% !TeX encoding = UTF-8

\documentclass[fontset=none]{exam-zh}
\usepackage{mathtools}
\usepackage{tcolorbox}
% 431 项目统一文档主题。
% 各根入口应先声明 \documentclass{exam-zh}，再输入本文件。

% 项目统一中文字体：避免依赖某台电脑上的 STSong/STHeiti 等系统字体。
\setCJKmainfont{SourceHanSerifCN-Regular.ttf}[
    Path = fonts/,
    BoldFont = SourceHanSerifCN-Bold.ttf,
    ItalicFont = AR-PL-KaitiM-GB.ttf,
    BoldItalicFont = AR-PL-KaitiM-GB.ttf,
    BoldItalicFeatures = {FakeBold = 2},
]
\setCJKsansfont{SourceHanSansCN-Regular.ttf}[
    Path = fonts/,
    BoldFont = SourceHanSansCN-Bold.ttf,
]
\setCJKmonofont{SourceHanSansCN-Regular.ttf}[
    Path = fonts/,
]

% fontset=none 时 CTeX 不创建这些常用字体族命令，显式补齐。
\setCJKfamilyfont{zhsong}{SourceHanSerifCN-Regular.ttf}[Path = fonts/]
\setCJKfamilyfont{zhhei}{SourceHanSansCN-Bold.ttf}[Path = fonts/]
\setCJKfamilyfont{zhkai}{AR-PL-KaitiM-GB.ttf}[Path = fonts/]
\setCJKfamilyfont{zhfs}{SourceHanSerifCN-Regular.ttf}[Path = fonts/]
\providecommand{\songti}{\CJKfamily{zhsong}}
\providecommand{\heiti}{\CJKfamily{zhhei}}
\providecommand{\kaishu}{\CJKfamily{zhkai}}
\providecommand{\fangsong}{\CJKfamily{zhfs}}


\usepackage{booktabs}
\usepackage{multirow}
\usepackage{threeparttable}
\usepackage{threeparttablex}
\usepackage{tikz}
\usetikzlibrary{calc}
\usetikzlibrary{arrows.meta}
\usetikzlibrary{decorations.pathmorphing}
\usetikzlibrary{decorations.pathreplacing}
\usetikzlibrary{positioning}

\tikzset{
    mindmap_root/.style={rectangle, rounded corners, draw=black, fill=gray!10, thick, align=center, minimum height=0.8cm, font=\bfseries\small},
    mindmap_level1/.style={rectangle, rounded corners, draw=blue!60, fill=blue!5, thin, align=center, minimum height=0.6cm, font=\small},
    mindmap_level2/.style={rectangle, draw=none, fill=none, align=center, font=\footnotesize, text width=3cm},
    arrow/.style={->, >=stealth, thick, gray}
}

\usepackage{tabularx}
\usepackage{xcolor}
\usepackage{makecell}
\usepackage{extarrows}
\usepackage{longtable}
\usepackage{float}
\usepackage{enumitem}
\usepackage{perpage}
\usepackage{scalerel}
\usepackage{stackengine}
\MakePerPage{footnote}
\stackMath
\tcbuselibrary{breakable}
\tcbset{breakable}

% 适用于较长估计变量名的宽帽号。
\newcommand{\reallywidehat}[1]{%
  \savestack{\tmpbox}{\stretchto{%
    \scaleto{%
      \scalerel*[\widthof{\ensuremath{#1}}]{\kern-.6pt\bigwedge\kern-.6pt}{%
        \rule[-\textheight/2]{1ex}{\textheight}%
      }%
    }{\textheight}%
  }{0.5ex}}%
  \stackon[1pt]{#1}{\tmpbox}%
}

% 列表全局设置
\setlist{nosep}
% 列表内部保持正文节奏；整个列表块与上下内容各保留半行间距。
\setlist[itemize]{
    leftmargin=*, labelsep=0.5em,
    itemsep=0pt, parsep=0pt,
    topsep=0.5\baselineskip, partopsep=0pt,
    after={}
}
\setlist[itemize,1]{label=$\bullet$}
\setlist[itemize,2]{label=\small$\bullet$}
\setlist[itemize,3]{label=\scriptsize$\bullet$}
\setlist[enumerate]{
    leftmargin=*, labelsep=0.5em,
    itemsep=0pt, parsep=0pt,
    topsep=0.5\baselineskip, partopsep=0pt,
    after={}
}
\setlist[description]{leftmargin=\parindent, labelindent=0pt}

% 章节公共样式。需要每节另起一页的资料在入口中单独覆盖 section/break。
\IfFileExists{ctex.sty}{
    \ctexset{
        section = {
            format = \Large\bfseries\raggedright,
            beforeskip = 1.5ex plus 0.5ex minus 0.2ex,
            afterskip = 1.5ex plus 0.2ex,
        },
        subsection = {
            format = \large\bfseries\raggedright,
            beforeskip = 1.0ex plus 0.5ex minus 0.2ex,
            afterskip = 1.0ex plus 0.2ex,
        },
        subsubsection = {
            format = \normalsize\bfseries\raggedright,
            beforeskip = 0.5ex plus 0.5ex minus 0.2ex,
            afterskip = 0.5ex plus 0.2ex,
        },
        paragraph = {
            runin = false,
            number = \arabic{paragraph}.,
            format = \normalsize\bfseries\raggedright,
            indent = 2em,
            beforeskip = 0.5ex plus 0.3ex minus 0.1ex,
            afterskip = 0.3ex plus 0.2ex,
        }
    }
}{}

% 按 chapter 计数器自动生成“第 X 讲”。
% exam-zh 在 \chapter 完成时推进 chapter，因此标题使用当前值加一。
\ExplSyntaxOn
\NewDocumentCommand{\lecture}{m}
  {
    \exp_args:Nx \chapter
      {第\zhnumber{\int_eval:n {\value{chapter} + 1}}讲：\exp_not:n {#1}}
  }
\ExplSyntaxOff

% 431 资料通用的 exam-zh 设置；可选参数为居中页眉，必选参数为右侧页眉。
\newcommand{\SetupFourThreeOneDocument}[2][]{
    \examsetup{
        page={
            size = a4paper,
            show-head = true,
            head-content = {
                \fancyhf{}
                \fancyhead[C]{#1}
                \fancyhead[EL]{\leftmark}
                \fancyhead[OR]{#2}
            },
            foot-content = {\hfill 第;页 / 共;页 \hfill}
        },
        style/student-version-cleanaux = true,
        fillin/type = paren,
        question = {
            show-points = auto,
            show-answer = true,
        },
        problem = {
            points-separate-par = false,
        },
        paren = {
            text-color = blue,
            show-paren = true,
            show-answer = true,
        },
        solution = {
            show-solution = show-stay,
            label-indentation = true,
            text-color = blue,
            blank-type = hide,
        },
    }

    \ExamPrintAnswerSet{
        question/show-answer = false,
        paren/show-answer = false,
        solution/show-solution = hide,
        paren/text-color = black,
        solution/text-color = black,
    }
}

\renewcommand{\chaptermark}[1]{\markboth{#1}{}}
\renewcommand{\headrulewidth}{1pt}

\newcommand{\true}{$\checkmark$}
\newcommand{\false}{$\times$}

\usepackage{hyperref}
\usepackage{tocloft}

\renewcommand{\contentsname}{\hspace*{\fill}目录\hspace*{\fill}}
\setlength{\cftbeforetoctitleskip}{2em}
\setlength{\cftaftertoctitleskip}{2em}
\setlength{\cftchapindent}{0em}
\setlength{\cftsecindent}{2.0em}
\setlength{\cftsecnumwidth}{2em}
\setlength{\cftsubsecindent}{4.0em}
\setlength{\cftsubsecnumwidth}{2.0em}
\setlength{\cftsubsubsecindent}{6.0em}
\setlength{\cftsubsubsecnumwidth}{2.0em}

\setlength{\intextsep}{1em}

\hypersetup{
    colorlinks = true,
    linkcolor = blue,
    urlcolor = blue,
    citecolor = blue,
    pdfborder = {0 0 0},
}

\newcommand{\kref}[1]{\ref{#1} \textsuperscript{(p.\pageref{#1})}}

\usepackage[group-separator={,}]{siunitx}

\tcbuselibrary{breakable}
\tcbset{breakable}

\SetupFourThreeOneDocument{金融经济学·熊和平·章末习题}

\title{金融经济学}
\subject{熊和平《金融经济学》本章小结与章末习题}
\author{Bilibili@园丁小区詹姆斯}

\examsetup{
    problem = {
        label = {},
        label-align = left,
        show-points = false,
    }
}

\newcommand{\problemno}[1]{\textbf{#1}\quad}

% 重要概念：与王江知识框一致，用条目列表而非逗号段。
\newcommand{\xhpconceptitems}[1]{%
\begin{tcolorbox}[knowledge,title=重要概念]
\begin{enumerate}[nosep]#1\end{enumerate}
\end{tcolorbox}}

\begin{document}

\FourThreeOneFrontMatter

\ifdefined\FourThreeOnePrintInterior
\else
\MakeFourThreeOneCover
    {431 FINANCE \quad/\quad FINANCIAL ECONOMICS}
    {金融经济学·熊和平·章末习题}
    {熊和平《金融经济学》本章小结与章末习题}
    {按教材章序整理本章小结、思考题与练习题（含参考解答）}
    {Bilibili@园丁小区詹姆斯}
    {PREFERENCE \quad EQUILIBRIUM \quad PRICING}

\null
\thispagestyle{empty}
\fi

\MakeFourThreeOneContents[2]

\cleardoublepage

\setcounter{chapter}{0}
% !TeX root = ../../金融经济学_熊和平_习题.tex
% 源：重拍 113800（印张 19）、113807（印张 20）；225493.pdf p.1

\zhchapter{基本框架及分析模式}

\begin{tcolorbox}[summary,title=本章小结]
金融经济学在分析问题时需要建立一定的模型，这种模型即是描述问题所涉及的基本框架。经济框架一般涉及三个方面：其一是经济环境；其二是证券市场；其三是经济人。

本书采用简单的单期框架，主要关注两个时点：期初投资和期末获得收益。经济环境还包括信息结构，即未来世界所处的状态及其分布。本书强调证券市场这一核心概念。经济人包括个人、厂商和政府；家庭金融把家庭从个人中分离出来单独研究。

一般均衡分析分两步：第一步，在偏好、禀赋和市场结构的约束下，经济人进行优化，确定消费或资产需求与价格之间的函数关系；第二步，利用市场出清条件（总禀赋等于总需求）联立求解均衡价格与配置。无套利分析法通过比较证券与组合的支付向量，寻找复制关系，从而对目标证券定价；其运用要求证券可被复制，并通常假设至少一位经济人具有不满足性、套利机会可被迅速利用（本书第六章展开具体方法）。

纯交换经济下证券供给外生给定，配置的可行性与 Pareto 最优（定义 1.2，约束 (1.15)--(1.16)）是评价资源配置的重要概念；引入生产后即为生产性经济，第十三章专门讨论。
\end{tcolorbox}

\xhpconceptitems{
    \item 经济环境，信息结构，事件，样本空间
    \item 投资组合，支付向量，支付空间，市场结构
    \item 一般均衡分析，套利分析，Lucas 树
    \item 风险证券，无风险证券，市场化，禀赋，消费计划，经济人
    \item 相对定价法，绝对定价法，市场出清，配置，Pareto 最优配置
}

% !TeX root = ../../金融经济学_熊和平_习题.tex

\begin{problem}
\problemno{1.1} 金融经济学是从哪几个方面来描述经济环境的？
\end{problem}
\begin{problem}
\problemno{1.2} 什么是纯交换经济？什么是生产性经济？两者的区别和联系是什么？
\end{problem}
\begin{problem}
\problemno{1.3} 金融经济学从哪些方面描述经济人？
\end{problem}
\begin{problem}
\problemno{1.4} 什么叫市场结构？
\end{problem}
\begin{problem}
\problemno{1.5} 什么是投资组合？投资组合一般有哪几种描述方式。
\end{problem}
\begin{problem}
\problemno{1.6} 什么是一般均衡分析方法？一般均衡分析方法一般包括哪几步？
\end{problem}
\begin{problem}
\problemno{1.7} 什么是无套利分析法？请简述该方法的基本原理。
\end{problem}


% !TeX root = ../../金融经济学_熊和平_习题.tex
% 源：重拍_20260929/IMG_20260929_113323.jpg（印张 42）

\zhchapter{期望效用理论}

\begin{tcolorbox}[summary,title=本章小结]
冯·诺依曼-摩根斯坦（Von Neumann--Morgenstern，1944）和萨维奇（Savage，1954）提出了期望效用理论，该理论是不确定条件下选择理论的主流理论，建立在公理化体系基础之上。该理论通常包括完备性、传递性、独立性、连续性和排序性公理等；若增加“存在最好彩票和最差彩票”的假设，可以较容易地证明效用函数的存在性，并且该效用函数表现为期望效用函数。期望效用函数具有正单调变换不变性。

两时点单期情形是本书的重点：单期情形下期望效用函数通常可以表示为二元效用函数的期望值；在状态独立、时间可加的假设下可以简化为由一元函数组成的期望效用函数。多于两时点的多期情形下同样可以将效用函数化简为易于处理的形式。

期望效用理论创立不久，心理学实验对该理论的若干公理提出了挑战，其中最著名的是阿莱斯悖论，这一悖论驱使经济学家提出了展望理论，并与其他理论一起促使行为金融的诞生和发展。行为金融的发展给期望效用理论提出了很多替代性理论，如一阶风险厌恶理论、状态依赖效用理论等；也有从纯粹技术角度对状态独立、时间可加效用函数进行拓展，比如递归效用函数。
\end{tcolorbox}

\xhpconceptitems{
    \item 不确定性，彩票，简单彩票，复合彩票，退化彩票
    \item 期望效用函数，圣彼得堡悖论，阿莱斯悖论
    \item 偏好逆转，框架效应，展望理论，消费习惯，公理化体系
}

% !TeX root = ../../金融经济学_熊和平_习题.tex

\begin{problem}
\problemno{2.1} 什么是不确定性？什么是可量化的不确定性？什么是不可量化的不确定性？
\end{problem}
\begin{solution}
不确定性是结果不能在决策时确定的情形。若能为各结果指定客观或主观概率，就是可量化不确定性（风险）；若无法合理指定概率，就是不可量化不确定性（模糊性）。
\end{solution}
\begin{problem}
\problemno{2.2} 为什么可以用彩票代替不确定性选择对象？
\end{problem}
\begin{solution}
彩票把每种结果及其概率明确列出，因此把复杂的不确定选择转化为对概率分布的比较。在期望效用公理成立时，任何不确定对象都可由与其结果—概率分布相同的彩票代表。
\end{solution}
\begin{problem}
\problemno{2.3} 什么是公理化体系？期望效用理论包含哪些公理？各个公理的具体含义是什么？
\end{problem}
\begin{solution}
公理化体系用一组行为公理刻画偏好，并由此推出表示定理。期望效用理论的核心公理包括完备性与传递性（偏好排序一致）、连续性（中间彩票可使偏好无差异）和独立性（共同混合不改变排序）；在这些条件下偏好可表示为 $U(L)=\sum_i\pi_i u(x_i)$。
\end{solution}
\begin{problem}
\problemno{2.4} 什么是效用函数？什么是期望效用函数？
\end{problem}
\begin{solution}
效用函数把消费或结果映射为实数，并保持偏好排序：$x\succeq y\Leftrightarrow u(x)\ge u(y)$。期望效用函数把各结果效用按概率加权：$U(L)=\sum_i\pi_i u(x_i)$。
\end{solution}
\begin{problem}
\problemno{2.5} 阿莱斯悖论挑战了期望效用理论的什么公理？
\end{problem}
\begin{solution}
阿莱斯悖论主要挑战独立性公理：受试者在两组选择中的共同结果相同，却因概率结构改变而改变偏好，违反共同概率混合应保持排序的要求。
\end{solution}
\begin{problem}
\problemno{2.6} 举例说明期望效用理论有哪些拓展。
\end{problem}
\begin{solution}
常见拓展包括累积前景理论（对概率作决策权重并以参考点定义收益损失）、秩依概率模型（按结果排序后加权）、模糊厌恶模型（对概率分布本身不确定）和非期望效用的加权效用模型。
\end{solution}
\begin{problem}
\problemno{2.7} 假设有两个彩票 $L$ 与 $L'$，以及效用函数 $V$ 与 $W$。若 $V(L)>V(L')$，是否必有 $W(L)>W(L')$？换言之，什么条件下 $W(L)=\sum_i \pi_i Q_i$ 也是(与 $V$ 表示相同偏好的)期望效用函数？% 原书 p.2 末题干截断，p.3 起为续问，待补扫核对完整设问。
\end{problem}
\begin{solution}
不一定。只有当 $W$ 与 $V$ 表示相同偏好时，才有 $W=f(V)$，其中 $f$ 为严格递增函数；若二者都是期望效用表示，则在相同结果空间和概率下，通常要求 $w(x)=a v(x)+b$ 且 $a>0$，于是 $W(L)=aV(L)+b$（或按归一化常数作等价变换），从而 $V(L)>V(L')\Leftrightarrow W(L)>W(L')$。
\end{solution}
\begin{problem}
\problemno{2.8} 资产 $A$ 具有确定性支付 \num{1500} 元；资产 $B$ 以 $0.8$ 的概率支付 \num{2000} 元、以 $0.2$ 的概率支付 \num{0} 元；资产 $C$ 以概率 $0.75$ 支付 \num{100} 元、以概率 $0.25$ 支付 \num{0} 元；资产 $D$ 以概率 $0.2$ 支付 \num{2000} 元、以概率 $0.8$ 支付 \num{0} 元。在期望效用框架下比较上述资产(按原书要求完成相应选择或排序)。
\end{problem}
\begin{solution}
仅凭支付和概率不能给出唯一排序，必须知道结果效用函数。期望支付分别为 $1500,1600,75,400$，但风险厌恶者可能偏好确定的 $A$ 胜过 $B$；风险中性时按期望支付排序为 $B>A>D>C$。
\end{solution}
\begin{problem}
\problemno{2.9} 若你具有对数效用函数，当前财富水平是 \num{5000} 元：
\begin{enumerate}[label=(\alph*)]
    \item 假如你面临一个 $50/50$ 的概率获得或损失 \num{1000} 元的风险，购买 \num{125} 元的保险可以完全消除该风险，你选择购买保险还是进行该项投机活动？
    \item 若在 (a) 中进行了投机而且亏损 \num{1000} 元，你的财富变为 \num{4000} 元。面对 (a) 中的保险是再次投机还是购买保险？
\end{enumerate}
\end{problem}
\begin{solution}
(a) 保险后的确定财富为 $5000-125=4875$。投机的期望对数效用为 $\frac12\ln6000+\frac12\ln4000=\frac12\ln(24\,000\,000)>\ln4875$，故选择投机。(b) 亏损后财富为 $4000$，再次投机的期望效用为 $\frac12\ln5000+\frac12\ln3000$，而买保险为 $\ln3875$；前者更大，仍选择投机。
\end{solution}
\begin{problem}
\problemno{2.10} 考虑下面的四种彩票：
\[
    L_A=(x,\pi;0,1-\pi),\quad L_B=(y,\pi';0,1-\pi'),\quad L_C=(x,\alpha\pi;0,1-\alpha\pi),\quad L_D=(y,\alpha\pi';0,1-\alpha\pi'),
\]
其中 $0<x<y$，$\pi<\pi'$，$\pi x<\pi' y$，并且 $\alpha\in(0,1)$。
\begin{enumerate}[label=(\alph*)]
    \item 若投资者在 $L_C$ 与 $L_A$ 之间选择 $L_C$，那么该投资者的偏好是否与冯·诺依曼-摩根斯坦的期望效用理论一致？为什么？
    \item 若投资者在 $L_A$ 与 $L_B$ 之间选择 $L_A$，在 $L_C$ 与 $L_D$ 之间选择 $L_D$，那么该投资者的偏好是否与冯·诺依曼-摩根斯坦的期望效用理论一致？为什么？
\end{enumerate}
\end{problem}
\begin{solution}
(a) 选择 $L_C$ 而非 $L_A$ 意味着在相同的中奖额下偏好更低中奖概率，违反不满足性，因而不符合期望效用理论。(b) 期望效用下 $L_A\succ L_B$ 意味 $\pi u(x)>(\pi')u(y)$；乘以 $\alpha$ 后应有 $L_C\succ L_D$，但题设选择 $L_D$，违反独立性公理，故不一致。
\end{solution}

% !TeX root = ../../金融经济学_熊和平_习题.tex
% 源：重拍_20260929/IMG_20260929_113635.jpg、113640.jpg（印张 62--63）

\zhchapter{风险厌恶及典型的效用函数}

\begin{tcolorbox}[summary,title=本章小结]
经济人对风险的选择由其风险态度描述，一般包括风险厌恶、风险中性和风险喜好，通常以是否愿意参与期望收益为零的公平博弈来判断：不愿参与者为风险厌恶，愿意参与者为风险喜好，无所谓者为风险中性。通常假设经济人为风险厌恶；风险厌恶对应凹效用函数，由概率论中的詹森不等式可见效用凹性与风险厌恶的关系；凸效用意味着风险喜好，线性效用意味着风险中性。马科维茨给出的风险溢价常作为风险厌恶的度量之一。

不同经济人的风险厌恶程度不同。阿罗—普瑞特用泰勒展开给出风险溢价的近似，并定义绝对风险厌恶系数与相对风险厌恶系数；确定性等价是使投资者无差异于随机财富的确定财富水平。文献中有两种确定性等价的常用定义，计算时略有差别。

两种风险厌恶系数虽来自近似，但在金融中有广泛应用。绝对风险厌恶为常数可定义 CARA 效用，相对风险厌恶为常数可定义 CRRA 效用；还可定义 HARA 及 IARA、DARA、IRRA、DRRA 等性质。CRRA 是 DRRA 的一个特例。
\end{tcolorbox}

\xhpconceptitems{
    \item 风险态度，风险厌恶，风险中性，公平博弈，詹森不等式
    \item 确定性等价，风险溢价，绝对风险厌恶系数，相对风险厌恶系数
    \item 风险容忍系数，谨慎系数，CARA，CRRA，HARA，风险厌恶比较
}

% !TeX root = ../../金融经济学_熊和平_习题.tex

\begin{problem}
\problemno{3.1} 什么是风险态度？一般包括哪几种？
\end{problem}
\begin{solution}
风险态度是经济人对随机财富的偏好，通常分为风险厌恶、风险中性和风险偏好。
\end{solution}
\begin{problem}
\problemno{3.2} 风险态度与效用函数的凹凸性之间具有什么样的关系？
\end{problem}
\begin{solution}
在边际效用为正的条件下，凹效用函数表示风险厌恶，线性效用表示风险中性，凸效用表示风险偏好。
\end{solution}
\begin{problem}
\problemno{3.3} 什么是詹森不等式？该不等式与风险厌恶有何关系？
\end{problem}
\begin{solution}
若 $u$ 为凹函数，则詹森不等式为 $u(E X)\ge E[u(X)]$。因此风险厌恶者偏好确定的期望财富，二者差额就是风险带来的效用损失。
\end{solution}
\begin{problem}
\problemno{3.4} 风险厌恶有哪几种度量方式？各种度量方法之间有何区别与联系？
\end{problem}
\begin{solution}
绝对风险厌恶系数为 $A(w)=-u''(w)/u'(w)$，以财富单位度量风险；相对风险厌恶系数为 $R(w)=-wu''(w)/u'(w)=wA(w)$，以财富比例度量风险。前者适合比较相同金额风险，后者适合比较相同比例风险。
\end{solution}
\begin{problem}
\problemno{3.5} 基于风险厌恶系数可以定义哪几种常见的效用函数？
\end{problem}
\begin{solution}
常见族包括绝对风险厌恶系数恒定的 CARA 指数效用、相对风险厌恶系数恒定的 CRRA 幂效用，以及二次效用和包含它们的 HARA 族。
\end{solution}
\begin{problem}
\problemno{3.6} 如何比较不同经济人之间的风险厌恶程度？
\end{problem}
\begin{solution}
在给定财富水平比较 $A(w)$ 或 $R(w)$；系数较大者更风险厌恶。若要作全局比较，应要求一个经济人的系数在所有相关财富水平都不小于另一个经济人。
\end{solution}

\begin{problem}
\problemno{3.7} 证明，对于 HARA 偏好：当 $\eta=0$ 时，HARA 效用函数变成了 CRRA 效用函数；当 $\gamma\to\infty$ 时，HARA 效用函数变成了 CARA 效用函数；当 $\gamma=-1$ 时，HARA 效用函数变成了二次效用函数。
\end{problem}
\begin{solution}
在 HARA 标准形式 $u(w)=\frac{(w+\eta)^{1-\gamma}}{1-\gamma}$($\gamma\neq1$)下，$\eta=0$ 时为 CRRA 幂效用；令 $\gamma\to\infty$ 并作适当标准化，极限为 $-\exp\{-aw\}$ 型 CARA 效用；$\gamma=-1$ 时指数为 $2$，得到二次效用函数。
\end{solution}
\begin{problem}
\problemno{3.8} 证明定理 3.1 中 (2) 与 (3) 等价。
\end{problem}
\begin{solution}
因 $u'(w)>0$，有 $A(w)=-u''(w)/u'(w)>0\Leftrightarrow -u''(w)>0\Leftrightarrow u''(w)<0$，所以两条件等价。
\end{solution}
\begin{problem}
\problemno{3.9} 用泰勒展开分别推导两种风险厌恶系数 ARA 和 RRA。
\end{problem}
\begin{solution}
对公平风险 $g$ 展开：$E u(w+g)=u(w)+u''(w)\operatorname{Var}(g)/2$。令 $u(w-\pi)=E u(w+g)$，得 $\pi\simeq A(w)\operatorname{Var}(g)/2$；除以财富后得到相对风险溢价 $\pi/w\simeq R(w)\operatorname{Var}(g)/(2w)$。
\end{solution}
\begin{problem}
\problemno{3.10} 已知指数效用函数 $u(w)=-\exp\{-aw\}$($a>0$)：
\begin{enumerate}[label=(\alph*)]
    \item 此函数是否表明边际效用为正且风险厌恶？
    \item 此函数是否意味着风险厌恶递减？
    \item 此效用函数是否意味相对风险厌恶为常数？
\end{enumerate}
\end{problem}
\begin{solution}
$u'(w)=ae^{-aw}>0$ 且 $u''(w)=-a^2e^{-aw}<0$，故边际效用为正且风险厌恶。$A(w)=a$ 为常数，$R(w)=aw$ 随财富线性增加，因此绝对风险厌恶恒定、相对风险厌恶不恒定。
\end{solution}
\begin{problem}
\problemno{3.11} 假设你具有 CRRA 偏好，相对风险厌恶系数为 $2$。如果你现在有 $50/50$ 的概率盈利或亏损 \num{1000} 元，同时规避风险的保险成本为 \num{500} 元，在什么财富水平下，投机活动与支付保险之间无差异？即确定性等价的财富是多少？
\end{problem}
\begin{solution}
取与 $R=2$ 等价的效用函数 $u(w)=-1/w$。无差异条件为
\[
u(w-500)=\frac12u(w-1000)+\frac12u(w+1000),
\]
整理无差异方程
\[
\frac1{w-500}=\frac{w}{w^2-1000^2}
\]
得 $w=2000$。因此确定性等价财富为 $w-500=1500$。
\end{solution}
\begin{problem}
\problemno{3.12} 经济人的初始财富为 $w$，而且他的绝对风险厌恶系数为常数 $a$，他需要承担风险 $\tilde{g}$。$\tilde{g}$ 是公平博弈，当 $\tilde{g}$ 具有如下分布时分别计算该经济人的确定性等价，并分析确定性等价与初始财富 $w$ 的关系：
\begin{enumerate}[label=(\alph*)]
    \item 取值为 $-6$ 和 $6$ 的二项分布；
    \item $[-c,c]$ 上的均匀分布；
    \item 均值为 $0$、方差为 $\sigma^2$ 的正态分布。
\end{enumerate}
\end{problem}
\begin{solution}
CARA 效用下
\[
CE=w-\frac1a\log E(e^{-a\tilde g}).
\]
二项分布 $\{-6,6\}$ 时 $CE=w-a^{-1}\log\cosh(6a)$；均匀分布 $[-c,c]$ 时 $CE=w-a^{-1}\log[\sinh(ac)/(ac)]$；正态分布时 $CE=w-a\sigma^2/2$。三种情形中风险成本均与 $w$ 无关。
\end{solution}
\begin{problem}
\problemno{3.13} 经济人的初始财富为 $w$，且绝对风险厌恶系数为 $1$，他需要承担风险 $w\tilde{g}$。$\tilde{g}$ 是公平博弈，当 $\tilde{g}$ 具有如下分布时分别计算该经济人的确定性等价，并分析确定性等价与初始财富 $w$ 的比值：
\begin{enumerate}[label=(\alph*)]
    \item 取值为 $-b$ 和 $b$ 的二项分布，其中 $0<b<1$；
    \item $[-c,c]$ 上的均匀分布，其中 $0<c<1$。
\end{enumerate}
\end{problem}
\begin{solution}
此时 $CE=w-\log E(e^{-w\tilde g})$。二项分布给 $CE=w-\log\cosh(wb)$，均匀分布给 $CE=w-\log[\sinh(wc)/(wc)]$。因此 $CE/w$ 分别取决于 $wb$、$wc$，不再是与财富无关的常数。
\end{solution}
\begin{problem}
\problemno{3.14} 经济人的初始财富为 $w$，且相对风险厌恶系数为 $2$，他需要承担风险 $\tilde{g}$，$\tilde{g}$ 取值为 $-b$ 和 $b$，$0<b<w$。有两种风险溢价的定义：
\[
    \pi_1(w)=\frac{E[u(w+\tilde{g})]}{u'(w)},\qquad \pi_2(w)=\frac{u(w)-E[u(w+\tilde{g})]}{u'(w)}.
\]
请问这两种方法计算的风险溢价相等吗？
\end{problem}
\begin{solution}
取 $u(w)=-1/w$。代入两种定义可得
\[
\pi_1=-\frac{w^3}{w^2-b^2},\qquad \pi_2=-\frac{wb^2}{w^2-b^2}.
\]
二者不相等；通常采用第二种定义，因为它对应使确定性财富下降到风险彩票期望效用的确定性金额。
\end{solution}
\begin{problem}
\problemno{3.15} 经济人的初始财富为 $w$，且是严格风险厌恶的。$\tilde{g}_1$ 和 $\tilde{g}_2$ 是两个独立的公平博弈，具有相同的分布，均为取值为 $-b$ 和 $b$ 的二项分布，其中 $0<b<w/2$。
\begin{enumerate}[label=(\alph*)]
    \item 假定他必须承担风险 $\tilde{g}_1$，则他的期望效用为 $V_1=E[u(w+\tilde{g}_1)]$。证明风险使他的情况恶化，即 $V_1=E[u(w+\tilde{g}_1)]<u(w)$。
    \item 假定他必须承担风险 $\tilde{g}=0.5(\tilde{g}_1+\tilde{g}_2)$，记在此情形下他的期望效用为 $V_2=E[u(w+\tilde{g})]$。请问 $V_2$ 是否总比 $V_1$ 高？
\end{enumerate}
\end{problem}
\begin{solution}
(a) 由严格詹森不等式及 $E\tilde g_1=0$，有 $V_1=E[u(w+\tilde g_1)]<u(w)$。(b) 平均风险 $\tilde g=(\tilde g_1+\tilde g_2)/2$ 的均值相同而分布更集中；严格凹性因此给出 $V_2>V_1$。
\end{solution}


% !TeX root = ../../金融经济学_熊和平_习题.tex
% 源：DocumentScanner_20260918_225493.pdf，扫描 p.6--7（印张约 85）

\zhchapter{Arrow-Debreu 经济}

\begin{tcolorbox}[knowledge,title=重要概念]
Arrow-Debreu 证券，状态价格，完全市场，金融财富，预算约束，市场出清，一般均衡，欧拉条件。
\end{tcolorbox}

% !TeX root = ../../金融经济学_熊和平_习题.tex

\xhpthoughts
\begin{problem}
\problemno{1} 什么是 Arrow-Debreu 证券市场？为什么要讨论 Arrow-Debreu 证券？
\end{problem}
\begin{problem}
\problemno{2} Arrow-Debreu 证券与现实生活中的什么金融产品类似？
\end{problem}
\begin{problem}
\problemno{3} 金融文献中在分析问题时经常假设市场是完全的，试说明完全市场的优点是什么？
\end{problem}
\begin{problem}
\problemno{4} Arrow-Debreu 证券市场的一般均衡分析分哪几步进行？
\end{problem}

\xhpexercises
\begin{problem}
\problemno{1} 考虑三状态经济中有三个证券，价格分别为 \num{5} 元、\num{5} 元和 \num{6} 元，期末的支付矩阵为
\[
    X=\begin{bmatrix}
        10 & 20 & 30\\
        10 & 10 & 10\\
        10 & 5 & 5
    \end{bmatrix}.
\]
\begin{enumerate}[label=(\arabic*)]
    \item 将三个证券表示为 A-D 证券的组合；
    \item 计算三个不同的状态证券价格。
\end{enumerate}
\end{problem}
\begin{problem}
\problemno{2} 给定下列支付矩阵，问 $k$ 为何值时市场是完全的：
\[
    X=\begin{bmatrix}
        3 & 2 & 3\\
        1 & k & 5\\
        2 & 1 & 7
    \end{bmatrix}.
\]
\end{problem}
\begin{problem}
\problemno{3} 市场上有三个证券，其价格和期末支付分别为
\[
    S=\begin{bmatrix}6\\4\\10\end{bmatrix},\qquad
    X=\begin{bmatrix}
        4 & 3 & 2\\
        2 & 7 & 3\\
        1 & 4 & 3
    \end{bmatrix}.
\]
\begin{enumerate}[label=(\arabic*)]
    \item 证明不存在套利机会；
    \item 确定状态价格向量；
    \item 确定消费计划 $(1,0,3)$ 的无套利价值；
    \item 确定以第一个证券为标的、执行价格为 $2$ 的看涨期权价格。
\end{enumerate}
\end{problem}
\begin{problem}
\problemno{4} 假定市场上有两种证券：证券 1 和证券 2，价格分别为 \num{13} 元和 \num{10} 元，两种证券期末的支付如下：
\begin{center}
\begin{tabular}{ccc}
    & 状态证券 1 的支付 & 状态证券 2 的支付\\
    \hline
    证券 1 & 11 & 0\\
    证券 2 & 2 & 20
\end{tabular}
\end{center}
假设经济人的效用函数为 $U(w_1,w_2)=w_1 w_2$，其中 $w_1,w_2$ 分别为期末不同状态下的财富水平。该经济人初始财富水平为 \num{420} 元。为了最大化效用，在期初构造投资组合时分别购买多少证券 1 和证券 2？
\end{problem}
\begin{problem}
\problemno{5} 假若期末只有一个状态发生，即期末是确定性的。一个消费者的两期消费的效用函数为 $u(c_0,c_1)=c_0 c_1^{1/2}$，其中 $c_0,c_1$ 分别为期初和期末的消费；该消费者的收入流为 $y_0=\num{2000}$ 元，$y_1=\num{1296}$ 元，且市场利率为 $8\%$。求该消费者的最优消费计划。
\end{problem}
\begin{problem}
\problemno{6} 继续考虑期末确定情形，经济人 1 在 0 期的禀赋为 $100$ 而在 1 期的禀赋为 $1$，其效用函数为 $u(c_0,c_1)=\ln c_0+\rho\ln c_1$；市场上只有一个无风险证券，无风险利率为 $r_f$。
\begin{enumerate}[label=(\arabic*)]
    \item 若 $\rho=0.5$，求经济人的投资需求；
    \item 若经济人 1 是唯一的参与者，请描述市场出清条件；
    \item 求均衡利率；
    \item 求均衡利率与时间偏好率 $\rho$ 的关系。
\end{enumerate}
\end{problem}
\begin{problem}
\problemno{7} 在第 4 题的基础上加入经济人 2，其效用函数与经济人 1 相同，禀赋为 $(1,100)$，即期初禀赋为 $1$，期末禀赋为 $100$。
\begin{enumerate}[label=(\arabic*)]
    \item 求经济人 2 的最优消费-投资策略；
    \item 求新的经济下的均衡利率，并比较单个经济人参与的经济下的均衡利率。
\end{enumerate}
\end{problem}
\begin{problem}
\problemno{8} 假设期末有三种等可能的状态 $a$、$b$ 和 $c$ 发生，市场是完全的，有两个经济人。经济人 1 的禀赋是 $100$ 个单位期初的消费品；经济人 2 的禀赋为期末的消费品，三个状态下分别为 $200$、$100$ 和 $50$。他们的效用函数相同，均为
\[
    u(c_0;c_a,c_b,c_c)=-\frac{1}{c_0}-\frac{1}{3}\Bigl(\frac{1}{c_a}+\frac{1}{c_b}+\frac{1}{c_c}\Bigr).
\]
\begin{enumerate}[label=(\arabic*)]
    \item 求解各自优化问题，得到每个人的最优消费-投资；
    \item 求均衡条件下三种状态的状态价格；
    \item 分析每个经济人的总消费与各自禀赋的关系。
\end{enumerate}
\end{problem}


% !TeX root = ../../金融经济学_熊和平_习题.tex
% 源：DocumentScanner_20260918_225493.pdf，扫描 p.8--9（印张约 108；p.9 含第五章练习题续，附录线性代数部分未收录）

\zhchapter{套利定价的理论基础}

% !TeX root = ../../金融经济学_熊和平_习题.tex

\begin{problem}
\problemno{5.1} 什么是套利？套利一般有哪几个类别？
\end{problem}
\begin{solution}
套利是一个组合 $\theta$，满足初始成本 $S^{\mathsf T}\theta\le0$，且期末支付 $X\theta\ge0$、至少一个状态严格为正。按来源可分为：利用同一支付不同价格的空间套利，利用不同到期价格关系的时间套利，利用复制关系的证券套利，以及利用统计定价偏离的统计套利。前三类是无套利定价理论的直接对象。
\end{solution}
\begin{problem}
\problemno{5.2} 什么是无套利原理？无套利原理成立的前提假设是什么？
\end{problem}
\begin{solution}
无套利原理要求不存在满足上述不等式的组合。其常用前提是：证券可按任意实数数量线性组合，市场价格对投资者给定，卖空和借贷在模型中可行，交易成本、税收和流动性限制忽略，且投资者严格偏好每个状态都更多的财富。若存在交易摩擦，应改用带摩擦的无套利条件。
\end{solution}
\begin{problem}
\problemno{5.3} 如何利用无套利原理对证券进行定价？
\end{problem}
\begin{solution}
设目标证券支付为 $Y$。若存在组合 $\theta$ 使 $X\theta=Y$，则复制组合成本 $S^{\mathsf T}\theta$ 必须等于目标价格 $P$；否则买入便宜的一方、卖出昂贵的一方即可在支付相同的情况下获得正初始收入。若只能找到 $X\theta\ge Y$，则 $S^{\mathsf T}\theta$ 给出目标价格的上界；反向支配关系给出下界。
\end{solution}
\begin{problem}
\problemno{5.4} 为什么要讨论证券市场的不同描述？
\end{problem}
\begin{solution}
支付矩阵 $X$ 直接描述每支证券在每个状态的支付，价格向量 $S$ 描述期初价值，组合空间 $\{X\theta\}$ 描述可实现支付，状态价格 $q$ 则把支付映射为价格 $S=X^{\mathsf T}q$。支付矩阵适合判断秩和完全性，状态价格适合定价和风险中性转换，组合空间适合直接检验套利。
\end{solution}
\begin{problem}
\problemno{5.5} 回忆线性代数中的秩的概念，它在证券市场的描述中起何作用？
\end{problem}
\begin{solution}
支付矩阵的秩是列空间维数，即独立可交易支付的数量。若有 $S$ 个状态，$\operatorname{rank}(X)=S$ 时每个状态支付都可复制，市场完全；若秩小于 $S$，至少有一个状态支付不能由证券组合生成。若证券数 $N$ 大于秩，则存在非零 $\theta$ 使 $X\theta=0$，证券中有冗余组合。
\end{solution}
\begin{problem}
\problemno{5.6} 冗余证券在套利中起什么作用？
\end{problem}
\begin{solution}
若 $X_j=X_{-j}\alpha$，证券 $j$ 的支付可由其他证券组合复制。无套利要求 $S_j=S_{-j}^{\mathsf T}\alpha$。价格高于右侧时卖出证券 $j$、买入复制组合；价格低于右侧时反向交易。两种情况下期末净支付均为零，期初都有正收入，因此冗余证券是检验一价律的直接工具。
\end{solution}
\begin{problem}
\problemno{5.7} 什么是一价律？一价律在衍生证券定价中是如何应用的？
\end{problem}
\begin{solution}
一价律是指 $X\theta=X\theta'$ 必有 $S^{\mathsf T}\theta=S^{\mathsf T}\theta'$。对衍生证券，先求组合 $\theta$ 使 $X\theta=Y$，则其无套利价格为 $P=S^{\mathsf T}\theta$。例如看涨期权可以由股票和无风险证券复制；若只能构造支付处处不低于 $Y$ 的组合，则复制成本是价格上界。
\end{solution}
\begin{problem}
\problemno{5.8} 什么是资产定价基本定理？理解其“基本”的含义。
\end{problem}
\begin{solution}
在有限状态模型中，资产定价基本定理表明以下命题等价：市场无套利；存在严格正的状态价格向量 $q$ 使 $S=X^{\mathsf T}q$；存在与客观概率等价的风险中性概率，使所有资产的贴现期望支付等于当前价格。“基本”是因为它把套利、线性状态价格和风险中性定价三种表述统一起来，是后续衍生证券定价的逻辑基础。
\end{solution}

\begin{problem}
\problemno{5.9} 证明定理 5.8 中的 (1)。
\end{problem}
\begin{solution}
若存在严格正的 $q$ 满足 $S=X^{\mathsf T}q$，设 $\theta$ 是套利组合。由价格关系，
\[
S^{\mathsf T}\theta=(X^{\mathsf T}q)^{\mathsf T}\theta=q^{\mathsf T}X\theta.
\]
套利定义给出 $S^{\mathsf T}\theta\le0$，而 $X\theta\ge0$ 且至少一个分量严格为正。因为每个 $q_s>0$，所以 $q^{\mathsf T}X\theta>0$，从而 $S^{\mathsf T}\theta>0$，与套利的初始成本条件矛盾。因此不存在套利。
\end{solution}
\begin{problem}
\problemno{5.10} 证券市场上有三个证券，其支付矩阵为
\[
    X=\begin{bmatrix}
        30 & 20 & 10\\
        20 & 15 & 0
    \end{bmatrix}.
\]
\begin{enumerate}[label=(\alph*)]
    \item 是否存在冗余证券？哪个证券是冗余的？请证明之。
    \item 若存在套利机会，请给出套利策略。
\end{enumerate}
\end{problem}
\begin{solution}
\begin{enumerate}[label=(\alph*)]
\item 由于只有两个状态而有三个证券，三列必线性相关。具体地，
\[
X_3=3X_1-4X_2
=3\begin{bmatrix}30\\20\end{bmatrix}
-4\begin{bmatrix}20\\15\end{bmatrix}
=\begin{bmatrix}10\\0\end{bmatrix}.
\]
因此证券 3 是冗余证券；它的支付可以由买入 3 单位证券 1、卖空 4 单位证券 2 完全复制。
\item 设三证券价格为 $S_1,S_2,S_3$。复制组合的成本为 $3S_1-4S_2$。若
\[
S_3>3S_1-4S_2,
\]
则卖出 1 单位证券 3，同时买入 3 单位证券 1、卖空 4 单位证券 2，期初获得正现金 $S_3-3S_1+4S_2$，期末支付恒为零，构成套利。若 $S_3<3S_1-4S_2$，则反向交易构成套利。因此无套利的必要条件是 $S_3=3S_1-4S_2$；在该等式成立时，仅由这组支付关系不能再构造价格套利。
\end{enumerate}
\end{solution}
\begin{problem}
\problemno{5.11} 若市场上三个证券的支付矩阵为
\[
    X=\begin{bmatrix}
        110 & 100 & 48\\
        110 & 50 & 40\\
        110 & 40 & 36
    \end{bmatrix},
\]
\begin{enumerate}[label=(\alph*)]
    \item 市场是否完全？请证明之。
    \item 若三个证券的价格分别为 $100$、$70$ 和 $40$，则是否存在套利机会？如果不存在套利机会，请计算无风险利率和 A-D 证券价格以及等价鞅测度。
    \item 若三个证券的价格分别为 $100$、$40$ 和 $70$，则是否存在套利机会？
    \item 上述三个复合证券分别是由怎样的 A-D 证券组合复合而成的？
\end{enumerate}
\end{problem}
\begin{solution}
\begin{enumerate}[label=(\alph*)]
\item
\[
\det X=13200\ne0,
\]
所以三列线性无关，且秩等于状态数 $3$，市场完全。
\item 令 $q=(q_1,q_2,q_3)^{\mathsf T}$。由 $X^{\mathsf T}q=(100,70,40)^{\mathsf T}$ 得
\[
110(q_1+q_2+q_3)=100,\quad100q_1+50q_2+40q_3=70,
\]
\[
48q_1+40q_2+36q_3=40.
\]
解得 $q=(17/33,3/11,4/33)^{\mathsf T}>0$，故不存在套利。总状态价格为 $8/11$，所以
\[
1+r_f=\frac{11}{8},\qquad r_f=\frac38.
\]
等价鞅概率为 $\widehat p=q/(8/11)=(17/24,3/8,1/6)^{\mathsf T}$。
\item 对价格 $(100,40,70)$ 解同一方程会得到负的状态价格分量。由资产定价基本定理，不存在严格正状态价格，故存在套利；套利组合可由该线性系统的负系数方向构造。
\item 三个证券分别是三列支付向量对应的状态证券组合：
\[
X_1=110A_1+110A_2+110A_3,quad
X_2=100A_1+50A_2+40A_3,
\]
\[
X_3=48A_1+40A_2+36A_3.
\]
\end{enumerate}
\end{solution}
\begin{problem}
% 原书 p.108：支付矩阵 OCR 残损，按三状态三资产整理，待补扫核对。
\problemno{5.12} 已知某个市场的支付矩阵为
\[
    X=\begin{bmatrix}
        1 & 0 & 0\\
        1 & 2 & 1\\
        1 & 3 & b
    \end{bmatrix},
\]
相应的价格向量为 $S=(6,31,9)^{\mathrm{T}}$。
\begin{enumerate}[label=(\alph*)]
    \item 证明存在冗余证券；
    \item $b$ 为何值时不存在套利机会？
    \item 不存在套利机会时计算状态价格、风险中性概率和利率。
\end{enumerate}
\end{problem}
\begin{solution}
\begin{enumerate}[label=(\alph*)]
\item 三个证券的支付列向量位于三维状态空间；是否冗余要看列之间的线性关系。由矩阵可知 $X_1=(1,1,1)^{\mathsf T}$、$X_2=(0,2,3)^{\mathsf T}$、$X_3=(0,1,b)^{\mathsf T}$。当 $b=3/2$ 时 $X_3=\frac12X_2$，证券 3 冗余。
\item 由价格方程
\[
q_1+q_2+q_3=6,\quad 2q_2+3q_3=31,\quad q_2+bq_3=9,
\]
可得
\[
q_3=\frac{13}{3-2b},\quad q_2=\frac{27-31b}{3-2b},\quad q_1=\frac{19b-22}{3-2b}.
\]
严格正状态价格同时要求 $b<27/31$（由 $q_2>0$）和 $b>22/19$（由 $q_1>0$），两者不可能同时成立。因此按当前整理出的支付矩阵和价格，不存在任何 $b$ 能使市场无套利；这也说明原书此处的 OCR 矩阵必须进一步核对。
\item 总状态价格恒为 $6$，故 $1/(1+r_f)=6$，即 $r_f=-5/6$。风险中性概率为 $\widehat p_s=q_s/6$，将上式分别除以 6 即得。
\end{enumerate}
\end{solution}
\begin{problem}
\problemno{5.13} 给定支付矩阵
\[
    X=\begin{bmatrix}
        2 & 3 & 2\\
        2 & 3 & 5\\
        5 & 4/3 & 3
    \end{bmatrix},
\]
和资产价格 $S^{\mathrm{T}}=(6,4,k)$：
\begin{enumerate}[label=(\alph*)]
    \item $K$ 为何值时不存在套利机会？
    \item 取 $K$ 为某特定值使得市场存在三类套利机会，分别给出相应的套利组合。
    \item 求以第一个资产为标的、以 $3$ 为执行价格的欧式看涨期权价格；求以第三个资产为标的、以 $2.5$ 为执行价格的欧式看涨期权价格。
    \item 若四个状态发生的概率相等，计算四个状态对应的随机贴现因子。
\end{enumerate}
\end{problem}
\begin{solution}
\begin{enumerate}[label=(\alph*)]
\item 令状态价格为 $q$，无套利条件为 $X^{\mathsf T}q=(6,4,k)^{\mathsf T}$ 且 $q>0$。先由前两列方程解出 $q_1,q_2,q_3$ 的关系，再将第三列代入得到 $k$ 的无套利取值（或区间）。
\item 若 $k$ 取使某个 $q_s=0$ 或 $q_s<0$ 的值，则分别沿该状态价格的负方向构造组合：当组合成本为零而支付非负且一项为正时为直接套利；当成本为负而支付为零时为融资套利；当成本和支付方向相反时买入低价组合、卖出高价组合。
\item 两个欧式看涨期权的支付向量分别为
\[
Y^{(1)}=(0,0,2)^{\mathsf T},\qquad Y^{(3)}=(0,3/2,1/2)^{\mathsf T}.
\]
完全市场下其价格分别为 $q^{\mathsf T}Y^{(1)}$ 与 $q^{\mathsf T}Y^{(3)}$。
\item 若四状态等概率 $p_s=1/4$，随机贴现因子为 $m_s=q_s/p_s=4q_s$。
\end{enumerate}
\end{solution}
\begin{problem}
\problemno{5.14} 证券市场上有两支股票，股票 A 和股票 B。股票 A 的期初价格为 \num{50} 元，期末两种状态 1 和 2 下的支付分别为 \num{60} 元和 \num{40} 元；股票 B 的期初价格为 \num{40} 元，期末两种状态下的支付分别为 \num{52} 元和 \num{28} 元。
\begin{enumerate}[label=(\alph*)]
    \item 给出价格向量和支付矩阵；
    \item 复制状态或有证券；
    \item 若状态 1 发生的概率为 $80\%$，则等权重组合的期望收益是多少？
\end{enumerate}
\end{problem}
\begin{solution}
\begin{enumerate}[label=(\alph*)]
\item 价格向量为 $S=(50,40)^{\mathsf T}$，支付矩阵为 $X=\begin{bmatrix}60&52\\40&28\end{bmatrix}$。
\item 复制状态 1 证券需解 $60\theta_A+52\theta_B=1$、$40\theta_A+28\theta_B=0$；解得 $\theta_A=-0.07,\theta_B=0.10$，其价格为 $50(-0.07)+40(0.10)=0.5$。复制状态 2 证券解得 $\theta_A=0.13,\theta_B=-0.15$，价格同样为 $0.5$。
\item 等权组合期初价格为 $45$，两状态支付分别为 $56$、$34$。状态 1 概率为 $0.8$ 时，期末期望支付为 $0.8\times56+0.2\times34=51.6$，故期望收益率为 $51.6/45-1=14.67\%$。
\end{enumerate}
\end{solution}


% !TeX root = ../../金融经济学_熊和平_习题.tex
% 源：DocumentScanner_20260918_225493.pdf，扫描 p.10--11（印张约 139--140）

\zhchapter{无套利原理在衍生证券定价中的应用}

\begin{tcolorbox}[summary,title=本章小结]
衍生证券一般包括远期、期货、期权、互换以及其他的混合证券。它们有着一些共同的特点，就是其价值依赖一些更基础的资产，我们通常称之为标的资产。衍生证券几乎都能利用无套利原理来进行定价；无套利原理在衍生证券定价中又可以形象地表示为“积木分析法”，即将一些基本的衍生证券依据其损益图看成是积木，把同一衍生证券的空头和多头看作不同的积木，将这些“积木”进行组合形成一些新的“积木”（新的损益图），从而对衍生证券进行定价或形成一些新的交易策略。

在所有衍生证券当中，期权非常具有代表性。依据执行时点的不同规定将期权分为欧式期权和美式期权，依据买卖权利的不同分为看涨期权和看跌期权。结合上一章的无套利原理以及由无套利原理导出的定理对不同期权的性质展开分析，得到了期权定价性质、期权的平价关系和美式期权是否提前执行的性质。

无套利定价在期权上的应用还体现在二叉树模型中的应用，结合状态价格使定价更简单和直观。同时，利用状态分离指数化证券构造期权还可以使市场完全化。
\end{tcolorbox}

\begin{tcolorbox}[knowledge,title=重要概念]
衍生证券，远期，期货，期权，积木分析法，实值期权，虚值期权，内在价值，时间价值，平价关系，美式期权，欧式期权，看涨期权，看跌期权，二叉树模型，风险对冲，状态指数化证券。
\end{tcolorbox}

% !TeX root = ../../金融经济学_熊和平_习题.tex

\begin{problem}
\problemno{6.1} 什么是衍生证券？衍生证券通常包括哪些种类？
\end{problem}
\begin{solution}
衍生证券是价值由标的资产、利率、汇率或指数等基础变量决定的金融合约。其支付依赖标的资产价格或状态，常见种类包括远期、期货、期权、互换，以及由这些合约组合而成的结构化产品。
\end{solution}
\begin{problem}
\problemno{6.2} 什么是积木分析法？它是如何体现无套利思想的？
\end{problem}
\begin{solution}
积木分析法把复杂衍生证券拆分为若干基本支付积木，例如零息债券、股票和状态证券，再分别定价并相加。若两种组合在每个状态支付完全相同，它们的价格必须相同；否则买入便宜组合、卖出昂贵组合即可套利。
\end{solution}
\begin{problem}
\problemno{6.3} 如何确定期权的价值？什么是内在价值？什么是时间价值？
\end{problem}
\begin{solution}
期权价值由复制或无套利上下界确定。看涨期权的内在价值为 $\max(S-K,0)$，看跌期权的内在价值为 $\max(K-S,0)$；时间价值是期权价格减去内在价值，反映到期前标的价格和波动率可能带来的额外收益。
\end{solution}
\begin{problem}
\problemno{6.4} 什么是状态指数化证券？为什么利用它能够使市场完全化？如何使市场完全化？
\end{problem}
\begin{solution}
状态指数化证券只在某一状态支付一个单位，在其他状态支付零。若有 $S$ 个状态并引入 $S$ 个相互独立的状态证券，支付矩阵就是单位矩阵，秩为 $S$，因此任意状态支付都可复制，市场完全化。实际操作是为每个不能复制的状态增加相应的状态证券。
\end{solution}
\begin{problem}
\problemno{6.5} 以二叉树模型为例，说明期权定价有哪些核心思想。
\end{problem}
\begin{solution}
核心思想包括：第一，用股票和无风险资产构造与期权相同的复制组合；第二，由一价律令期权价格等于复制成本；第三，用风险中性概率计算贴现期望支付；第四，利用向后递推从到期节点逐层求出各时点价格。风险中性概率只是定价工具，不必等于真实上涨概率。
\end{solution}
\begin{problem}
\problemno{6.6} 分别说明美式期权在支付红利和不支付红利两种情形下是否提前执行，说明理由。
\end{problem}
\begin{solution}
不支付红利的股票看涨期权通常不应提前执行：持有期权保留了上行选择权，提前执行会失去时间价值。支付红利时，除息会使股票价格下降，若红利足够大，临近除息日提前执行可能最优。美式看跌期权即使不支付红利也可能提前执行，因为立即获得执行价与股票价值之差可能超过继续持有的时间价值。
\end{solution}
\begin{problem}
\problemno{6.7} 以期权价格的上下界为例说明这些等量关系和不等量关系成立的前提条件。
\end{problem}
\begin{solution}
无套利下欧式看涨期权满足 $\max(S_0-Ke^{-rT},0)\le c\le S_0$，欧式看跌期权满足 $\max(Ke^{-rT}-S_0,0)\le p\le Ke^{-rT}$。等式需要复制或退化条件，例如到期立即执行时内在价值等于价格；不等式依赖无套利、可借贷、无交易成本和标的可交易等前提。支付红利时上下界要用标的资产净现值替换 $S_0$。
\end{solution}

\begin{problem}
\problemno{6.8} 证明期权性质 1 至性质 6。
\end{problem}
\begin{solution}
各性质均由逐状态支付比较得到。例如看涨期权支付 $(S_T-K)^+$ 随 $S_T$ 单调增加、随 $K$ 单调减少，故 $c$ 随标的价格增加、随执行价格减少；期限延长扩大了可选择的执行时点，故美式期权价值不减；无套利给出 $c\ge\max(S_0-Ke^{-rT},0)$ 和 $p\ge\max(Ke^{-rT}-S_0,0)$。若任一关系违反，则买入支付较高而价格较低的组合、卖出支付较低而价格较高的组合即可构造套利。
\end{solution}
\begin{problem}
\problemno{6.9} 假设 $c_1$、$c_2$ 和 $c_3$ 分别表示执行价格分别为 $X_1$、$X_2$ 和 $X_3$ 的欧式看涨期权的价格，其中 $X_1<X_2<X_3$，并且 $X_3-X_2=X_2-X_1$。三个期权的标的资产相同、到期日相同。证明：$2c_2\le c_1+c_3$。
\end{problem}
\begin{solution}
对任意到期价格 $S_T$，函数 $c(K)=(S_T-K)^+$ 是执行价格 $K$ 的凸函数。由于 $X_2=(X_1+X_3)/2$，有
\[
2(S_T-X_2)^+\le(S_T-X_1)^++(S_T-X_3)^+.
\]
对状态价格加权并求和，得到 $2c_2\le c_1+c_3$。若不成立，则买入右侧较便宜的一方并卖出左侧较昂贵的一方即可套利。
\end{solution}
\begin{problem}
\problemno{6.10} 蝶形看涨价差是由标的资产相同、到期日相同而执行价格不相同的欧式看涨期权形成的期权组合：(1) 买入一份执行价格为 $K-\delta$ 的看涨期权和一份执行价格为 $K+\delta$ 的看涨期权；(2) 卖出 $2$ 份执行价格为 $K$ 的看涨期权。请画出蝶形价差的损益图，并证明当 $K$ 给定而 $\delta$ 不同时，$\delta$ 越大，蝶形价差的价格越高。
\end{problem}
\begin{solution}
到期支付为
\[
B(S_T)=(S_T-K+\delta)^+-2(S_T-K)^++(S_T-K-\delta)^+.
\]
因此 $S_T\le K-\delta$ 或 $S_T\ge K+\delta$ 时支付为零；在 $K-\delta<S_T\le K$ 时支付为 $S_T-K+\delta$，在 $K<S_T<K+\delta$ 时支付为 $K+\delta-S_T$，图形是顶点为 $(K,\delta)$ 的三角形。由凸性不等式，$c(K-\delta)+c(K+\delta)-2c(K)\ge0$；增大 $\delta$ 扩大非负支付区域，且逐状态支付不减，故无套利价格随 $\delta$ 增大而上升。
\end{solution}
\begin{problem}
\problemno{6.11} 考虑一个二叉树模型，股票 A 的期初价格为 \num{10} 元，上涨和下跌等概率，且 $u=20\%$，$d=-20\%$，以 A 为标的欧式看涨期权，无风险利率为 $10\%$。
\begin{enumerate}[label=(\alph*)]
    \item 画出股票 A 和执行价格为 \num{9} 元的欧式看涨期权的二叉树图形，计算期权在期初的价格；
    \item 假如得到新的利好消息：股票上涨的概率大于 $0.5$，股票期初价格仍为 \num{10} 元。请问你愿意购买看涨期权还是看跌期权？为什么？
\end{enumerate}
\end{problem}
\begin{solution}
\begin{enumerate}[label=(\alph*)]
\item 一期股票价格为 $12$ 或 $8$，执行价为 $9$ 的看涨期权支付为 $3$ 或 $0$。风险中性概率满足 $1.2p^*+0.8(1-p^*)=1.1$，故 $p^*=3/4$，期权初始价格为 $c_0=(3/4)3/1.1=45/22\approx2.045$。
\item 利好消息提高上涨的主观概率时，看涨期权的主观期望收益增加；若投资者相信该消息，应选择看涨期权。风险中性概率用于无套利定价，不等于真实概率。
\end{enumerate}
\end{solution}
\begin{problem}
\problemno{6.12} 在第 6.11 题的基础上考虑三时点两期的二叉树模型，上涨和下跌率保持不变，无风险利率也不变。
\begin{enumerate}[label=(\alph*)]
    \item 画出标的股票和欧式看涨期权的二叉树图形，并计算各节点的期权价格；
    \item 若期权为美式看涨期权，画出看涨期权的二叉树图形，计算各节点的期权价格。
\end{enumerate}
\end{problem}
\begin{solution}
取风险中性概率 $p^*=1/2$。四条路径的期末支付依次为 $7,1,0,0$，因为对应平均价格为 $12,6,3,1.5$。因此期初无套利价格为
\[
c_0=\frac{(7+1)/4}{1.25^2}=1.28。
\]
\end{solution}
\begin{problem}
\problemno{6.13} 同样考虑三时点(0、1、2 时点)两期二叉树模型，标的资产期初价格为 $S_0=4$ 元，$u=2$，$d=1/2$，无风险利率为 $25\%$。以该股票为标的资产的亚式看涨期权，期权最终价值为 $c_2=\max\bigl[(S_1+S_2)/2-5,0\bigr]$，即执行价格为 \num{5} 元，以股票在时点 1 和 2 处的价格 $S_1$、$S_2$ 的算术平均值来结算。请计算期权的期初无套利价格。
\end{problem}
\begin{solution}
取 $p^*=1/2$，四条路径的 $(S_1,S_2)$ 为 $(8,16),(8,4),(2,4),(2,1)$，平均价格为 $12,6,3,1.5$，期末支付为 $7,1,0,0$。因每条路径概率为 $1/4$，
\[
c_0=\frac{(7+1)/4}{(1.25)^2}=1.28。
\]
\end{solution}
\begin{problem}
\problemno{6.14} 假定期末有 $4$ 个不同的状态发生，假设股票 A 的支付向量为 $(1,2,3,4)^{\mathrm{T}}$，请你以股票 A 为基础构造欧式看涨期权使市场完全化。用记号表示期权和股票的价格，用这些符号表示出状态价格。
\end{problem}
\begin{solution}
选取执行价格 $1,2,3$ 的三个欧式看涨期权，连同股票 A 形成四个证券。支付向量为
\[
A=(1,2,3,4),\ C_1=(0,1,2,3),\ C_2=(0,0,1,2),\ C_3=(0,0,0,1)。
\]
其支付矩阵为上三角矩阵且行列式为 $1$，所以市场完全。若股票和期权价格分别为 $S,c_1,c_2,c_3$，状态价格满足
\[
q_4=c_3,\quad q_3=c_2-2c_3,\quad q_2=c_1-2c_2+c_3,\quad q_1=S-2c_1+c_2。
\]
\end{solution}


% !TeX root = ../../金融经济学_熊和平_习题.tex
% 源：重拍_20260929/IMG_20260929_113128.jpg（印张 160）、IMG_20260929_113355.jpg（印张 159）

\zhchapter{基于消费的投资组合理论}

\begin{tcolorbox}[summary,title=本章小结]
马科维茨的投资组合理论虽然是独立地提出来的，但是依然可以从最一般的情形找到理论依据。为了更好地理解这个理论，我们回到最基本的情形：在两时点单期情形下，投资者在自身禀赋约束下通过资本市场的交换重新配置资源，从而使其效用极大。为了追求其“一生”效用的极大化，资源需要在两时点重新配置，因此投资者问题归结为消费-组合问题。一方面，我们研究了消费-组合问题最优解的存在性；另一方面，在状态有限、市场完全下我们将问题化简为单纯的投资组合问题。

为了分析组合问题最优解的性质，我们首先分析了两资产情形下最优解的性质，得到了“最优风险资产买卖原则”：当且仅当风险资产的超额收益为正时，最优的策略是购买若干风险资产；当且仅当风险资产的超额收益为负时，最优的策略是卖空若干风险资产；当且仅当风险资产的超额收益为零时，最优的策略是全部持有无风险资产。

由于“最优风险资产买卖原则”不太符合直觉，我们在特定情形下分析了风险资产的选择问题。小风险情形下风险资产的最优投资数量近似地与风险资产超额收益成正比例，与风险成反比例，与绝对风险厌恶系数成反比例，这一结论在“CARA 加正态分布”情形是准确地成立的。进一步在“CRRA 加对数正态分布”情形下讨论得出最优风险资产的权重与风险资产超额收益成正比例、与风险成反比例、与相对风险厌恶系数成反比例。这些结论具有明确的经济含义并形成我们分析投资组合问题的“基准”。两资产情形下的部分结论可以推广到多资产情形时，发现除了少数情形下可以得到显示解，大多数情形下问题变得比较复杂。特别需要强调的是“CARA 加正态分布”情形可以推广到多元情形，三个及以上的风险资产满足该假设条件时也能得到显示解，因此很多文献喜欢采用 CARA 偏好。
\end{tcolorbox}

\xhpconceptitems{
    \item 消费-组合问题，风险资产买卖原则
    \item 小风险情形，投资需求弹性
}

% !TeX root = ../../金融经济学_熊和平_习题.tex

\begin{problem}
\problemno{7.1} 我们在分析消费-组合问题时为什么要转化为纯组合问题？这样对经济人的优化问题的求解是否产生影响？
\end{problem}
\begin{problem}
\problemno{7.2} 在分析小风险情形后为什么还要讨论 CARA 偏好下的情形？
\end{problem}
\begin{problem}
\problemno{7.3} CARA 偏好下一般能够得出显式解，为什么我们通常不假设 CARA 偏好？
\end{problem}
\begin{problem}
\problemno{7.4} HARA 偏好下的组合问题与 CARA 和 CRRA 偏好下的组合问题有何异同？
\end{problem}
\begin{problem}
\problemno{7.5} 两资产情形下的组合问题的结论哪些能够推广到多资产情形？
\end{problem}

\begin{problem}
\problemno{7.6} 两资产情形下(一个风险资产和一个无风险资产)，假设经济人具有对数效用函数，并且风险资产的收益率以等概率取无风险利率的 $4$ 倍和负的无风险利率。求投资者对风险资产的投资在其初始财富中所占的比例 $A/W_0$。
\end{problem}
\begin{problem}
\problemno{7.7} 假设投资者具有 CARA 偏好，绝对风险厌恶系数 $a>0$。投资者的期末财富服从正态分布，则投资者的确定性等价财富水平是多少？
\end{problem}
\begin{problem}
\problemno{7.8} 同第 1 题一样，考虑两资产情形：无风险资产收益率为 $r_f$，风险资产收益率为 $\tilde{r}$，其均值和方差分别为 $\bar{r}$ 和 $\sigma^2$：
\begin{enumerate}[label=(\alph*)]
    \item 在二次效用假设下求最优投资组合；并讨论最优组合与无风险利率 $r_f$、风险利率的期望值 $\bar{r}$ 以及方差 $\sigma^2$、偏好系数 $a$ 的关系，给出其经济学解释。此时效用函数为
    \[
        U=E\bigl[\tilde{W}_1-\tfrac12 a\tilde{W}_1^2\bigr].
    \]
    \item 在 CARA 偏好下重复 (1) 中的计算和分析。此时效用函数为
    \[
        U=E\bigl[-e^{-a\tilde{W}_1}\bigr].
    \]
\end{enumerate}
\end{problem}
\begin{problem}
\problemno{7.9} 证明定理 7.8。
\end{problem}
\begin{problem}
\problemno{7.10} 两资产情形下期末财富为 $W=w+a\tilde{x}$(如公式 (7.3) 所示)，$R(w)$ 为相对风险厌恶系数，$a(w)$ 为最优投资选择，即
\[
    a(w)=\operatorname*{argmax}_{a} E\bigl[u(w+a\tilde{x})\bigr],
\]
满足条件 $w>0$，$E[\tilde{x}]>0$，$u'(\cdot)>0$，$u''(\cdot)\le 0$。证明以下命题：
\begin{enumerate}[label=(\alph*)]
    \item 若 $\lambda>1$，且对于给定的效用函数 $u$ 有 $R'(w)\ge 0$。定义 $u_1(w)\equiv u(\lambda w)$，$u_2(w)=u(w)$，证明 $R_1(w)\ge R_2(w)$。
    \item 给定以上结论，我们有 $a_1(w)\le a_2(w)$，其中 $a_1(w)$、$a_2(w)$ 分别是 $u_1(\cdot)$、$u_2(\cdot)$ 对应的最优投资选择。
    \item 定义 $\hat{a}(w)=a(w)/w$，则对于 $w_1\le w_2$ 有 $\hat{a}(w_1)\ge \hat{a}(w_2)$(提示：令 $\lambda=w_2/w_1$)。
\end{enumerate}
\end{problem}


% !TeX root = ../../金融经济学_熊和平_习题.tex
% 源：微信图片_51253_3585、51254_3586（教材印张 182--183）

\zhchapter{马科维茨投资组合理论}

\begin{tcolorbox}[summary,title=本章小结]
马科维茨最早建立了投资组合理论，他用投资组合收益的期望值度量组合的收益，用方差或标准差度量组合的风险，从而提出了均值-方差标准。托宾则证明了均值-方差标准的适用条件：如果投资者的效用函数是二次函数或者风险资产服从联合正态分布，则均值方差标准有效。

在均值方差标准下，从两资产到多资产情形可以得到投资组合形成的可行集和有效集，即所有投资组合的收益与风险在均值-方差平面形成的集合。在这些可行组合中收益相同且风险最低的组合称为组合前沿；组合前沿上风险相同而收益最高的组合为有效组合。通常在均值-方差平面内的组合前沿为抛物线，或者在均值-标准差平面内为双曲线的右半支。引入无风险资产后的组合前沿是过无风险资产并与组合前沿相切的直线，均衡条件下形成了资本市场线。

在进行投资组合分析时，有马科维茨的直观分析，也有莫顿的严格数学分析。莫顿的分析给出了组合前沿组合、最小方差组合以及零协方差组合的计算公式。依照莫顿的结果可知：任何两个异于最小方差组合的前沿组合可以生成整个组合前沿；任何几个前沿组合或有效组合的凸组合依然是前沿组合或有效组合。
\end{tcolorbox}

\xhpconceptitems{
    \item 收益，收益率，相对收益，投资组合，风险，均值-方差标准
    \item 系统性风险，非系统性风险，特质风险
    \item 组合前沿，前沿组合，有效组合，有效集，可行集，资本市场线
    \item 最小方差组合，零协方差组合，协方差矩阵
}

% !TeX root = ../../金融经济学_熊和平_习题.tex

\begin{problem}
\problemno{8.1} 什么是相对收益？什么是收益率？两者之间有何关系？
\end{problem}
\begin{solution}
相对收益是期末支付与期初价格之比，记为 $R_i=S_{i1}/S_{i0}$；收益率是相对收益减去 1，即 $r_i=R_i-1=(S_{i1}-S_{i0})/S_{i0}$。相对收益包含本金，收益率只表示本金的百分比变化。
\end{solution}
\begin{problem}
\problemno{8.2} 如何计算投资组合的收益率？
\end{problem}
\begin{solution}
若组合权重为 $w_i$ 且 $\sum_iw_i=1$，则组合相对收益为 $R_p=\sum_iw_iR_i$，收益率为 $r_p=\sum_iw_ir_i$。权重按期初市值计算；若有现金流，应把分红和利息计入期末支付。
\end{solution}
\begin{problem}
\problemno{8.3} 什么是均值-方差标准？什么条件下该标准成立？
\end{problem}
\begin{solution}
均值—方差标准只用 $E(r_p)$ 和 $\operatorname{Var}(r_p)$ 比较组合：在期望收益相同的组合中选方差小者，在方差相同的组合中选期望收益大者。它在收益联合正态（或更一般的椭圆分布）时对所有二次或凹效用投资者成立；小风险二阶近似下也可作为近似标准。
\end{solution}
\begin{problem}
\problemno{8.4} 两风险资产情形下，什么时候组合前沿为曲线？什么情形下组合前沿是折线？
\end{problem}
\begin{solution}
若两资产收益不完全相关，协方差矩阵正定，组合方差关于权重是严格凸二次函数，前沿为平滑曲线。若相关系数为 $1$ 或 $-1$，一个资产可由另一个线性复制，风险—收益边界退化为由直线段组成的折线；允许卖空时还可能出现一条直线。
\end{solution}
\begin{problem}
\problemno{8.5} 什么是系统性风险？什么是非系统性风险？其中哪一种风险可以被分散化？如何分散？
\end{problem}
\begin{solution}
系统性风险由宏观因素和共同因子引起，影响大多数资产，不能通过简单分散消除；非系统性风险是单个企业或行业特有的，可通过持有相关性较低的大量资产分散。因子模型中总风险可写为因子风险与特质风险之和，组合扩大后特质风险趋近于零。
\end{solution}
\begin{problem}
\problemno{8.6} 在 $N$ 个风险资产中引入无风险资产后，对投资组合的可行集和有效前沿有何影响？
\end{problem}
\begin{solution}
无风险资产收益为 $r_f$、方差为零。任一风险组合与无风险资产的混合都位于均值—标准差图上的直线 $E(r)=r_f+\frac{E(r_p)-r_f}{\sigma_p}\sigma$。有效前沿变为从无风险点出发、与原风险资产有效边界相切的资本配置线；允许借贷时直线可延伸到切点两侧。
\end{solution}

\begin{problem}
\problemno{8.7}（与王江第 8 章 8.1 题相似。）假设 $N$ 个风险资产，收益率为 $\tilde{r}_i$($i=1,2,\ldots,N$)，它们是独立同分布的。对于一个具有不满足性和凹性效用函数的投资者，计算他的最优投资组合，并求均值-方差标准下的最优组合。
\end{problem}
\begin{solution}
独立同分布意味着每个资产的期望收益相同、方差相同且协方差为零。对称性和凹性使最优组合给每个资产相同权重 $w_i=1/N$；否则把权重从低边际效用资产转向高边际效用资产可提高效用。均值—方差标准下，组合期望收益不依赖权重，而方差为 $\sigma^2\sum_iw_i^2$，在 $\sum_iw_i=1$ 下由 Cauchy 不等式在 $w_i=1/N$ 处最小，方差为 $\sigma^2/N$。
\end{solution}
\begin{problem}
\problemno{8.8} 已知市场上有三个证券，其协方差矩阵为
\[
    V=\begin{bmatrix}
        24 & -10 & 9\\
        -10 & 75 & 3\\
        9 & 3 & 12
    \end{bmatrix}.
\]
\begin{enumerate}[label=(\alph*)]
    \item 计算等权重投资组合的方差。
    \item 计算如下两个投资组合的方差：组合 1：$(0.1,0.8,0.1)^{\mathrm{T}}$；组合 2：$(1.25,-0.1,-0.15)^{\mathrm{T}}$。
\end{enumerate}
\end{problem}
\begin{solution}
\begin{enumerate}[label=(\alph*)]
\item 等权重 $w=(1/3,1/3,1/3)^{\mathsf T}$，组合方差为 $w^{\mathsf T}Vw=115/9\approx12.778$。
\item 对 $w_1=(0.1,0.8,0.1)^{\mathsf T}$，$w_1^{\mathsf T}Vw_1=47.42$；对 $w_2=(1.25,-0.1,-0.15)^{\mathsf T}$，$w_2^{\mathsf T}Vw_2=37.735$。负权重表示卖空相应证券。
\end{enumerate}
\end{solution}
\begin{problem}
\problemno{8.9} 假设市场上有两项资产。资产 1 收益率的期望值为 $0.3$，方差为 $5.9$；资产 2 收益率的期望值为 $0.08$，方差为 $0.05$；两资产的协方差为 $-0.01$。
\begin{enumerate}[label=(\alph*)]
    \item 求由上述两个资产形成的可行集；
    \item 求最小方差集的组合向量。
\end{enumerate}
\end{problem}
\begin{solution}
设资产 1 权重为 $w$，资产 2 权重为 $1-w$。可行集参数式为
\[
\mu(w)=0.08+0.22w,\quad \sigma^2(w)=5.9w^2+0.05(1-w)^2-0.02w(1-w).
\]
最小方差权重由 $w^*=\frac{\sigma_2^2-\sigma_{12}}{\sigma_1^2+\sigma_2^2-2\sigma_{12}}=\frac{0.06}{5.97}\approx0.01005$，故组合向量约为 $(0.01005,0.98995)^{\mathsf T}$。
\end{solution}
\begin{problem}
\problemno{8.10} 三种资产的期望收益率向量和协方差矩阵为
\[
    ER=\begin{bmatrix}10.1\%\\7.8\%\\5\%\end{bmatrix},\qquad
    V=\begin{bmatrix}
        2.1 & 0.6 & 0\\
        0.6 & 0.9 & 0\\
        0 & 0 & 0
    \end{bmatrix}.
\]
\begin{enumerate}[label=(\alph*)]
    \item 三种资产中哪一个是无风险资产？为什么？
    \item 计算组合 $(0.5,0.5,0)$ 和组合 $(0.4,0.4,0.2)$ 的期望收益和标准差。
\end{enumerate}
\end{problem}
\begin{solution}
\begin{enumerate}[label=(\alph*)]
\item 第三个资产方差为零，因此它是无风险资产；其收益率为 $5\%$。
\item 组合 $(0.5,0.5,0)$ 的期望收益为 $8.95\%$，方差为 $0.25(2.1+0.9+2\times0.6)=1.05$，标准差为 $\sqrt{1.05}\approx1.025$（按题中百分比单位）。组合 $(0.4,0.4,0.2)$ 的期望收益为 $8.12\%$，方差为 $0.16(2.1+0.9+1.2)=0.672$，标准差约为 $0.820$。
\end{enumerate}
\end{solution}
\begin{problem}
\problemno{8.11} 证明定理 8.6，定理 8.7 和性质 8.2。
\end{problem}
\begin{solution}
组合方差为 $\sigma_p^2=w^{\mathsf T}Vw$，在 $\mathbf1^{\mathsf T}w=1$ 下使用拉格朗日乘子，一阶条件 $2Vw-\lambda\mathbf1=0$ 给出 $w\propto V^{-1}\mathbf1$，这证明最小方差组合的表达式。加入期望收益约束 $\mu^{\mathsf T}w=\bar\mu$ 后，解线性方程组得到均值—方差有效前沿；消去权重得到关于 $\bar\mu$ 的二次方程，故前沿为双曲线。若加入无风险资产，所有组合位于无风险点与切点组合的直线上，这就是资本配置线性质。
\end{solution}


% !TeX root = ../../金融经济学_熊和平_习题.tex
% 源：113071.pdf 第 3--4 页（印张 205--206）

\zhchapter{资本资产定价模型}

\begin{tcolorbox}[summary,title=本章小结]
如果在实证结果中表现出对股票收益具有显著的预测能力，我们都称之为异象。在过去的几十年间实证发现了大量的异象，例如一类称为日历异象的效应，这些效应包括：周末效应（weekend effect），一月效应（the January effect），月度效应（the monthly effect），假日效应（holiday effect）和日内效应（intraday effect）。

资本资产定价模型简称 CAPM，是建立在大量假设基础之上的理论；可以将假设分为关于主观和客观两个方面的假设。在主观假设下得到均值-方差（标准差）平面上的无差异曲线，在客观假设下可以得到同样平面内的组合前沿，两者结合可以得到单个投资者的最优组合选择。当所有投资者具有同质信念，同时资本市场具有 $N$ 个风险资产和一个无风险资产时，市场均衡使得过无风险资产并与 $N$ 个风险资产组合前沿相切的切点组合与市场投资组合重合，从而形成均衡条件下的资本市场线和证券市场线。

所谓资本市场线（CML）是指均衡条件下 $N+1$ 个证券组成的所有有效组合所形成的直线，换言之，直线上的每个点代表的组合是有效组合；反之，每一个有效组合必定位于 CML 直线上。所谓证券市场线（SML）是指均衡条件下所有证券或组合的收益与风险在均值-贝塔平面内形成的直线；均衡条件下凡是偏离 SML 的点所代表的证券或组合，其定价就出现错误。SML 给出了单个证券或证券组合的定价方程，它表明任何证券或证券组合的收益与其风险（贝塔系数）成直线关系。传统的实证方法包括截面回归和时间序列回归；实证结果与理论模型的不完全一致推动了资本市场定价理论的发展和完善。
\end{tcolorbox}

\xhpconceptitems{
    \item 同质信念，无差异曲线，分离定理
    \item 资本市场线（CML），证券市场线（SML），证券特征线，贝塔系数
    \item 零贝塔 CAPM，Fama 和 MacBeth 截面回归
    \item 金融异象，罗尔批判，Fama-French 三因子模型
}

% !TeX root = ../../金融经济学_熊和平_习题.tex

\begin{problem}
\problemno{9.1} 什么是同质信念？为什么要强调同质信念？
\end{problem}
\begin{solution}
同质信念是指所有投资者对各资产未来收益的概率分布、期望和协方差有相同判断。它使所有人面对同一个均值—方差机会集，从而共同选择同一个切点风险组合；个人差异只体现在无风险资产与该组合之间的配置比例。没有同质信念，市场组合未必就是每个人的切点组合，标准 CAPM 的分离结论会失效。
\end{solution}
\begin{problem}
\problemno{9.2} 我们将假设分为关于主观的假设和关于客观的假设，请问两部分假设各自与此前哪些章节的内容相关？
\end{problem}
\begin{solution}
主观假设涉及投资者的偏好、风险厌恶、期望效用和均值—方差选择，对应前面的效用函数、风险厌恶和投资组合理论；客观假设涉及资产支付、市场结构、交易制度、无风险借贷和市场出清，对应 Arrow--Debreu 经济、无套利定价和市场均衡部分。
\end{solution}
\begin{problem}
\problemno{9.3} 为什么说 CAPM 是一般均衡定价？
\end{problem}
\begin{solution}
CAPM 不是对单个投资者的局部最优条件，而是把所有投资者的最优组合、资产供给和市场出清同时联立。均衡中每个人持有无风险资产和同一个市场组合，价格调整使总需求等于资产总供给，最终得到 $E(R_i)=R_f+\beta_i[E(R_M)-R_f]$，所以它是一般均衡定价关系。
\end{solution}
\begin{problem}
\problemno{9.4} 什么是分离定理？分离定理中的“分离”是指什么与什么的分离？该定理在实践中所隐含的意义是什么？
\end{problem}
\begin{solution}
分离定理指出，在均值—方差条件下，所有投资者都持有同一个切点风险组合，只在无风险资产与该组合之间选择不同权重。“分离”是生产（或风险组合的确定）与个人消费偏好的分离：先确定共同风险组合，再由风险厌恶决定配置比例。实践中可由专业机构提供市场组合，个人只需按风险承受能力调整无风险资产比例。
\end{solution}
\begin{problem}
\problemno{9.5} 什么是资本市场线(CML)？其含义是什么。
\end{problem}
\begin{solution}
CML 是均值—标准差图中从无风险点 $(R_f,0)$ 出发、与风险资产有效前沿相切的直线：
\[
E(R_p)=R_f+\frac{E(R_M)-R_f}{\sigma_M}\sigma_p.
\]
它只适用于有效的有效组合，斜率是单位总风险的风险溢价，即市场组合的夏普比率。
\end{solution}
\begin{problem}
\problemno{9.6} 什么是证券市场线(SML)？CML、SML 以及证券特征线彼此有何区别？
\end{problem}
\begin{solution}
SML 为 $E(R_i)=R_f+\beta_i[E(R_M)-R_f]$，横轴是系统性风险 $\beta$，适用于单个证券和任意组合。CML 横轴是总风险标准差，只描述有效组合。证券特征线是时间序列回归 $R_{it}-R_{ft}=\alpha_i+\beta_i(R_{Mt}-R_{ft})+\varepsilon_{it}$，用于估计单个证券的贝塔和特质误差，并非均衡定价线。
\end{solution}
\begin{problem}
\problemno{9.7} 什么是零贝塔 CAPM？
\end{problem}
\begin{solution}
当不存在可按无风险利率无限借贷的资产时，可用与市场组合协方差为零的零贝塔组合替代无风险资产。零贝塔 CAPM 写为 $E(R_i)=E(R_Z)+\beta_i[E(R_M)-E(R_Z)]$，其中 $\operatorname{Cov}(R_Z,R_M)=0$，截距是零贝塔组合的期望收益。
\end{solution}
\begin{problem}
\problemno{9.8} 简单列举资本资产定价模型实证的经典方法。
\end{problem}
\begin{solution}
经典方法包括：时间序列回归检验证券特征线和零截距限制；横截面回归检验平均收益与贝塔的线性关系及市场风险溢价；Fama--MacBeth 两步回归；按贝塔排序形成组合后比较组合收益；以及联合检验市场组合代理、贝塔稳定性和残差风险是否定价。
\end{solution}
\begin{problem}
\problemno{9.9} 在 BJS 实证分析中计算各股票的贝塔值 $\beta_i$ 采用的是对式 (9.21)
\[
    R_{it}-R_{ft}=\alpha_i+\beta_i(R_{Mt}-R_{ft})+\varepsilon_{it}
\]
的回归，那么与利用贝塔公式计算的贝塔值有何异同？
\end{problem}
\begin{solution}
式 (9.21) 即证券特征线的时间序列回归 $R_{it}-R_{ft}=\alpha_i+\beta_i(R_{Mt}-R_{ft})+\varepsilon_{it}$，回归斜率
\[
\hat\beta_i=\frac{\widehat{\operatorname{Cov}}(R_i-R_f,R_M-R_f)}{\widehat{\operatorname{Var}}(R_M-R_f)}。
\]
与贝塔定义 $\beta_i=\operatorname{Cov}(R_i,R_M)/\operatorname{Var}(R_M)$ 在总体意义上等价（常用超额收益形式）；样本回归给出的是其估计量，并附带标准误与显著性检验。若用原始收益直接算协方差，在 $R_f$ 随时间变化时与回归结果可能不一致。
\end{solution}
\begin{problem}
\problemno{9.10} 什么是金融异象？金融异象对金融理论的发展有何作用？
\end{problem}
\begin{solution}
金融异象是不能由标准 CAPM 的贝塔—收益关系充分解释的系统性经验规律，如规模效应、价值效应、动量效应和低波动效应。它们促使研究者检验市场组合代理、风险因子和投资者理性假设，发展多因子模型、行为金融和条件资产定价模型；异象本身并不自动证明存在无风险套利，因为交易成本、风险暴露和数据挖掘也可能解释表面偏离。
\end{solution}

\begin{problem}
\problemno{9.11} 如果市场上无风险利率为 $5\%$，市场投资组合的收益率为 $15\%$，标准差为 $20\%$，则期望收益率为 $20\%$ 的证券其贝塔系数为多少？其标准差是多少？它与市场投资组合的相关系数是多少？如果另外一个股票的期望收益率为 $25\%$，方差为 $52\%$，则该股票的系统性风险和非系统性风险各是多少？
\end{problem}
\begin{solution}
\begin{enumerate}[label=(\alph*)]
\item 由 SML，$0.20=0.05+\beta(0.15-0.05)$，故 $\beta=1.5$。
\item 若该证券位于资本市场线，其标准差为 $\beta\sigma_M=1.5\times0.20=0.30$，并与市场完全正相关，相关系数为 $1$。仅由 SML 本身不能推出标准差；这里使用其为有效组合的条件。
\item 另一股票的 $\beta=(0.25-0.05)/0.10=2$。市场方差为 $0.20^2=0.04$，系统性方差为 $\beta^2\sigma_M^2=4\times0.04=0.16$；总方差为 $0.52$，故非系统性方差为 $0.52-0.16=0.36$。
\end{enumerate}
\end{solution}
\begin{problem}
\problemno{9.12} 资产 $X$ 与 $Y$ 的期望收益和协方差矩阵为
\[
    ER=\begin{bmatrix}0.2\\0.1\end{bmatrix},\qquad
    V=\begin{bmatrix}0.01&0\\0&0.0064\end{bmatrix}.
\]
\begin{enumerate}[label=(\alph*)]
    \item 市场投资组合的权重为 $X$ 和 $Y$ 各占 $50\%$，求零贝塔组合的期望收益。
    \item 求全局最小方差组合的组合权重。
    \item 求全局最小方差组合与零贝塔组合的协方差。
    \item 给出证券市场线的方程。
\end{enumerate}
\end{problem}
\begin{solution}
\begin{enumerate}[label=(\alph*)]
\item 市场组合为 $(1/2,1/2)$。资产 $X,Y$ 与市场组合的协方差分别为 $0.005$ 和 $0.0032$。令零贝塔组合权重为 $(z,1-z)$，则 $0.005z+0.0032(1-z)=0$，得 $z=-16/9$，所以其期望收益为 $0.2(-16/9)+0.1(25/9)=-7/90$。
\item 全局最小方差权重为
\[
w_X=\frac{0.0064}{0.01+0.0064}=\frac{16}{41},\qquad w_Y=\frac{25}{41}.
\]
\item 零贝塔组合的协方差为零是相对于市场组合而言；与全局最小方差组合的协方差为 $w_{\rm GMV}^{\mathsf T}Vw_Z=7/1792\approx0.00391$。
\item 零贝塔 CAPM 的证券市场线为
\[
E(R_i)=-\frac7{90}+\beta_i\left(\frac{3}{20}+\frac7{90}\right),
\]
其中 $\beta_i=\operatorname{Cov}(R_i,R_M)/\operatorname{Var}(R_M)$。
\end{enumerate}
\end{solution}
\begin{problem}
\problemno{9.13} 假设均值-方差机会集由两个资产 A 和 B 构成，其期望收益和方差-协方差矩阵为
\[
    ER=\begin{bmatrix}30\%\\20\%\end{bmatrix},\qquad
    V=\begin{bmatrix}0.0081&0\\0&0.0025\end{bmatrix}.
\]
\begin{enumerate}[label=(\alph*)]
    \item 假设投资者 I 选择的“市场投资组合”为 $(0.75,0.25)$，投资者 J 选择不同的“市场投资组合”为 $(0.5,0.5)$，在此背景下每个投资者计算的资产 A 的贝塔值是多少？
    \item 在此条件下两投资者对资产 A 要求的收益率各是多少？
    \item 对于每一个投资者，计算他们认定的零贝塔组合和证券市场线的方程。
\end{enumerate}
\end{problem}
\begin{solution}
记 $R_M^{(k)}=w_A^{(k)}R_A+w_B^{(k)}R_B$，$\beta_A^{(k)}=\operatorname{Cov}(R_A,R_M^{(k)})/\operatorname{Var}(R_M^{(k)})$。
\begin{enumerate}[label=(\alph*)]
\item 投资者 I：$E(R_M)=0.75\times30\%+0.25\times20\%=27.5\%$，$\operatorname{Var}(R_M)=0.75^2\times0.0081+0.25^2\times0.0025=0.0047125$，$\operatorname{Cov}(R_A,R_M)=0.75\times0.0081=0.006075$，故 $\beta_A^{(I)}=0.006075/0.0047125\approx1.29$。投资者 J：$E(R_M)=25\%$，$\operatorname{Var}(R_M)=0.00265$，$\operatorname{Cov}(R_A,R_M)=0.00405$，故 $\beta_A^{(J)}=0.00405/0.00265\approx1.53$。
\item 在各自以所选组合为“市场”的 CAPM 下，$E(R_A)=R_f+\beta_A^{(k)}[E(R_M^{(k)})-R_f]$。故 I 要求 $R_f+1.29\times(27.5\%-R_f)$，J 要求 $R_f+1.53\times(25\%-R_f)$；两式一般不等，说明“市场组合”认定不同会得到不同的定价线（除非 $R_f$ 取特殊值使两者同时等于 $30\%$）。
\item 零贝塔组合 $w_Z$ 满足 $\mathbf1^{\mathsf T}w_Z=1$ 且 $\operatorname{Cov}(R_Z,R_M^{(k)})=w_Z^{\mathsf T}Vw_M^{(k)}=0$。I 得 $w_Z^{(I)}\approx(-0.115,\,1.115)^{\mathsf T}$，$E(R_Z^{(I)})\approx18.85\%$，SML 为 $E(R)=18.85\%+\beta\,(27.5\%-18.85\%)$。J 得 $w_Z^{(J)}\approx(-0.446,\,1.446)^{\mathsf T}$，$E(R_Z^{(J)})\approx15.54\%$，SML 为 $E(R)=15.54\%+\beta\,(25\%-15.54\%)$。
\end{enumerate}
\end{solution}
\begin{problem}
\problemno{9.14} 某证券的市价为 \num{40} 元，该证券的期望收益率等于 $13\%$，无风险利率等于 $7\%$，市场风险溢价(市场投资组合期望收益率与无风险利率之差)为 $8\%$。如果该证券在未来的支付保持不变，但它与市场投资组合的协方差变为原来的两倍，则证券的价格应该为多少？
\end{problem}
\begin{solution}
原来的贝塔为 $\beta=(0.13-0.07)/0.08=0.75$。协方差加倍而市场方差不变，新的贝塔为 $1.5$，故新的期望收益率为 $0.07+1.5\times0.08=0.19$。未来支付不变时，价格与贴现率满足 $P\propto1/(1+R)$；原价格为 $40$、原收益率为 $13\%$，故未来支付的价值为 $40(1.13)$，新价格为
\[
P'=\frac{40(1.13)}{1.19}\approx37.98\text{ 元}。
\]
\end{solution}


% !TeX root = ../../金融经济学_熊和平_习题.tex
% 源：113071.pdf p.5--6（印张 225--226）；小结见重拍 113742（印张 224）

\zhchapter{因子模型与套利定价理论}

\begin{tcolorbox}[summary,title=本章小结]
CAPM 表明，均衡条件下任何证券或证券组合的期望收益与市场投资组合的期望收益呈线性关系，由此形成证券市场线（SML）。SML 表明任何证券的收益由市场投资组合这一单一因子决定，推而广之，证券收益由某个单一因子决定的模型即是单因子模型。因子模型在求证券期望收益和方差及协方差时大大简化。当单一因子模型不能很好地解释证券收益时，需要将单因子模型推广到两因子乃至多因子模型；多因子模型同样能简化证券收益和风险的确定。

在多因子模型基础上，罗斯提出了套利定价模型（APT），套利定价建立在较少的假设基础上。为推导定价方程，我们先研究了精确因子模型，然后得到极限套利和套利定价模型的定价方程。

套利定价理论或多因子定价模型非常利于实证分析，因子模型的实证分析本质上是资本资产定价模型实证分析的延伸，经典的实证分析包括 Fama--French 三因子模型、五因子模型和其他的多因子模型。
\end{tcolorbox}

\xhpconceptitems{
    \item 市场模型，单因子模型，多因子模型，因子暴露，因子载荷
    \item 因子风险，非因子风险，APT，精确因子，因子溢价
    \item 风险价格，极限套利，因子构建
}

% !TeX root = ../../金融经济学_熊和平_习题.tex

\begin{problem}
\problemno{10.1} 什么是市场模型？市场模型下证券的收益和风险如何确定？
\end{problem}
\begin{problem}
\problemno{10.2} 什么是多因子模型？多因子模型下证券的收益如何确定？证券的风险如何分解？
\end{problem}
\begin{problem}
\problemno{10.3} APT 有哪些假设？
\end{problem}
\begin{problem}
\problemno{10.4} 简述 APT 与 CAPM 的联系与区别。
\end{problem}
\begin{problem}
\problemno{10.5} Fama-French 三因子模型如何构造决定证券收益的因子？
\end{problem}
\begin{problem}
\problemno{10.6} 如果 CAPM 和 APT 都是正确的资产定价理论，那么与市场投资组合负相关的因子的风险溢价必定是负的，为什么？从数值和直观上进行解释。
\end{problem}
\begin{problem}
\problemno{10.7} 有学者认为市场投资组合是不能被测定的，因此 CAPM 是不可检测的。也有学者认为 APT 既未明确因子的数目也未明确因子的内容，从而也是不可检测的。如果这些观点都是正确的，这意味着这些理论毫无价值吗？解释之。
\end{problem}
\begin{problem}
\problemno{10.8} 虽然 APT 未明确有关因子的内容，多数套利定价理论的实证研究集中在某些类型的因子上，这些因子的共同特征是什么？
\end{problem}

\begin{problem}
\problemno{10.9} 在单因子模型下，市场上有两个风险资产 A 和 B，两证券的风险载荷分别为 $0.20$ 和 $3.5$；两者的特质风险(方差)分别为 $49\%$ 和 $100\%$。当两资产组成的投资组合为 $(40\%,60\%)$ 时分别回答下列各问题：
\begin{enumerate}[label=(\alph*)]
    \item 如果单因子的标准差为 $15\%$，则投资组合的因子风险是多少？
    \item 组合的非因子风险和组合的总风险各是多少？
    \item 如果无风险资产、A 和 B 组成的投资组合为 $(10\%,36\%,54\%)$，分别回答上述两问。
\end{enumerate}
\end{problem}
\begin{problem}
\problemno{10.10} 在两因子模型下，市场上有两个证券 A 和 B，它们对因子 1 的敏感性分别为 $1.5$ 和 $0.7$，对因子 2 的敏感性分别为 $2.6$ 和 $1.2$，两证券的特质风险用方差表示分别为 $25\%$ 和 $16\%$，因子 1 和因子 2 的标准差分别为 $20\%$ 和 $15\%$，因子间协方差为 $2.25\%$。计算两个证券的标准差和它们之间的协方差。
\end{problem}
\begin{problem}
\problemno{10.11} 单因子模型下，假设无风险利率为 $2\%$，一个具有单位敏感性的投资组合的期望收益率为 $4.5\%$。假设有两个风险证券 A 和 B 的风险载荷分别为 $4.0$ 和 $2.6$，根据套利定价理论，求这两种证券组成组合 $(30\%,70\%)$ 的均衡期望收益率。
\end{problem}
\begin{problem}
\problemno{10.12} 假设证券收益由两因子模型生成。两个风险证券和无风险资产对每个因子的敏感性及每种证券的期望收益率见下表。
\begin{center}
\begin{tabular}{cccc}
    证券 & $b_{i1}$ & $b_{i2}$ & 期望收益率(\%)\\
    \hline
    A & 0.50 & 0.80 & 5.2\\
    B & 1.50 & 1.40 & 7.6\\
    C & 0.00 & 0.00 & 2\\
\end{tabular}
\end{center}
\begin{enumerate}[label=(\alph*)]
    \item 如果你有 \num{100} 元用于投资，卖空 \num{50} 元证券 B 买入 \num{150} 元证券 A，你的投资组合对两个因子的敏感性是多少？
    \item 如果你以无风险利率借入 \num{100} 元，并将该资金与原来的钱合并按 (1) 中的比例投资于资产 A 和 B，则该组合对两个因子的敏感性各为多少？求该组合的期望收益率。
    \item 因子 2 的期望风险溢价是多少？
\end{enumerate}
\end{problem}
\begin{problem}
\problemno{10.13} 假设 CAPM 成立，证券的收益由单因子模型生成。给定如下信息：$\sigma_A^2=4.00\%$，$b_A=0.70$，$b_B=1.10$，$\operatorname{cov}(F,r_M)=3.70\%$。
\begin{enumerate}[label=(\alph*)]
    \item 计算证券 A 和 B 的贝塔系数；
    \item 均衡期望收益率各是多少？
\end{enumerate}
\end{problem}
\begin{problem}
\problemno{10.14} 假设 CAPM 成立，证券的收益由两因子模型生成。给定如下信息：$\sigma_A^2=3.24\%$，$b_{A1}=0.80$，$b_{B1}=1.00$，$\operatorname{cov}(F_1,r_M)=1.56\%$；$b_{A2}=1.10$，$b_{B2}=0.70$，$\operatorname{cov}(F_2,r_M)=5.00\%$。在上述条件下回答第 5 题中的两问。
\end{problem}
\begin{problem}
\problemno{10.15} 假设资产 A、B、C 的收益服从精确两因子模型：
\[
    r_A=a_A+b_{A1}F_1+b_{A2}F_2,\quad
    r_B=a_B+b_{B1}F_1+b_{B2}F_2,\quad
    r_C=a_C+b_{C1}F_1+b_{C2}F_2,
\]
其中两个因子 $F_i$($i=1,2$)是均值为 $0$、方差为 $1$ 的因子，$r_f$ 为无风险利率。无套利条件下 $a$、$b$ 应该满足什么条件？
\end{problem}


% !TeX root = ../../金融经济学_熊和平_习题.tex
% 源：113071.pdf p.7--9（印张约 252--253）

\zhchapter{完全市场下的资源配置与资产定价}

\begin{tcolorbox}[summary,title=本章小结]
完全市场下经济人的决策问题可以从经济人的禀赋入手进行分析；经济人的禀赋可能会呈多种形式，比如只有实物禀赋、实物加证券禀赋和既有实物又有证券禀赋，实际上三种禀赋都是可以互相转化的，因此可以在最简单的情形下展开分析。进一步，完全市场下，如果市场无摩擦，经济人在含 $N$ 个证券市场下的最优决策与 A-D 框架下的最优决策是等价的。

完全市场下的自由竞争均衡与 A-D 框架下的一般均衡是相同的，包含两个步骤：一是各自优化，二是市场出清。均衡条件下最优消费由各自的效用函数和状态价格确定；个人在不同状态下的消费水平由相应状态单位概率的状态价格确定，即每个状态下状态价格与该状态的概率之比决定该状态下的消费水平，这种比也称为随机贴现因子或状态价格密度。

在讨论多人参与的经济环境里，资源配置的好坏标准是 Pareto 有效。Pareto 最优配置等价于一个中央计划者的优化问题。福利经济学第一和第二定理在自由竞争问题和中央计划者问题之间建立起了密切的关系。在福利经济学定理的基础上，提出了风险分担规则的概念，并且可以证明线性风险容忍是线性分享规则的充要条件。此外，还可以建立代表性经济人的框架。
\end{tcolorbox}

\xhpconceptitems{
    \item 状态价格密度，资产定价的核，随机贴现因子（SDF）
    \item Pareto 最优，福利权重，福利经济学定理，Pareto 最优分享规则
    \item 谨慎系数，代表性经济人
}

% !TeX root = ../../金融经济学_熊和平_习题.tex

\begin{problem}
\problemno{11.1} 什么是资源配置？什么是可行配置？什么是 Pareto 有效配置？
\end{problem}
\begin{solution}
资源配置是把各期各状态的消费品和生产要素分配给所有经济人。配置可行是指每期每状态各人的消费总和不超过经济可用资源；Pareto 有效配置是可行配置中不存在使至少一人严格变好而不使任何人变坏的另一配置。
\end{solution}
\begin{problem}
\problemno{11.2} 什么是福利经济学第一定理？什么是福利经济学第二定理？
\end{problem}
\begin{solution}
第一定理：在局部非饱和、价格-taking 和市场出清条件下，竞争均衡配置是 Pareto 有效的。第二定理：在凸偏好、凸技术和适当连续性条件下，任一 Pareto 有效配置都可通过先进行一次性财富转移、再让市场竞争实现。第一定理说明市场配置的效率，第二定理说明公平与效率可分离。
\end{solution}
\begin{problem}
\problemno{11.3} 什么是风险分担规则？什么条件下风险分担规则是线性的？
\end{problem}
\begin{solution}
风险分担规则 $f_i(C)$ 给出总消费或总财富为 $C$ 时经济人 $i$ 在每种状态获得的消费。Pareto 条件为 $\lambda_i u_i'(f_i(C))$ 在经济人之间相等。由此，若所有经济人的风险容忍函数在最优配置上相同（特别是同类 CARA 或适当同类偏好），则 $f_i'(C)$ 为常数，风险分担规则为线性函数。
\end{solution}
\begin{problem}
\problemno{11.4} 什么是代表性经济人？代表性经济人的效用函数与单个经济人的效用函数之间有什么联系？
\end{problem}
\begin{solution}
代表性经济人是一个虚拟经济人，其偏好能对所有个人偏好加总，并在总资源约束下产生与分散经济人相同的总需求和资产价格。给定 Pareto 权重 $\lambda_i$，其状态效用为 $u_0(C)=\sum_i\lambda_i u_i(f_i(C))$，其中 $f_i$ 是 Pareto 风险分担规则；它不是简单把个人效用直接相加，而是先在每个总消费水平上进行最优分配。
\end{solution}
\begin{problem}
\problemno{11.5} 如何度量代表性经济人的风险厌恶系数？代表性经济人的风险容忍与单个经济人之间的风险容忍有什么关系？
\end{problem}
\begin{solution}
代表性经济人的绝对风险厌恶为 $A_0(C)=-u_0''(C)/u_0'(C)$，风险容忍为 $T_0(C)=1/A_0(C)=-u_0'(C)/u_0''(C)$。在 Pareto 最优分配下，风险容忍具有可加性：
\[
T_0(C)=\sum_iT_i(f_i(C))。
\]
因此总风险容忍等于个人风险容忍之和，代表性经济人的绝对风险厌恶通常低于任何单个经济人。
\end{solution}
\begin{problem}
\problemno{11.6} 什么是状态价格密度？
\end{problem}
\begin{solution}
状态价格密度是把某状态单位支付折算为当前价格的随机变量。若状态 $s$ 的客观概率为 $p_s$、状态价格为 $q_s$，则状态价格密度为 $m_s=q_s/p_s$。任一资产价格均可写为 $P=E[mX]$；它把状态价格向量与随机贴现因子统一起来。
\end{solution}
\begin{problem}
\problemno{11.7} 假设市场是完全的，在 1 期经济有 $3$ 个等可能的状态 $a$、$b$ 和 $c$；经济中有两个参与者 1 和 2，他们的禀赋分别为
\[
    e_1=\begin{bmatrix}100\\0\\0\end{bmatrix},\qquad
    e_2=\begin{bmatrix}0\\100\\50\end{bmatrix}.
\]
他们的时间偏好都为 $1$，而且都具有 CRRA 效用函数，参与者 1 的相对风险厌恶系数为 $1$，参与者 2 的相对风险厌恶系数为 $2$。
\begin{enumerate}[label=(\alph*)]
    \item 求解每一参与者的最优消费计划；
    \item 求解均衡状态价格；
    \item 求期末支付向量为 $(200,100,50)^{\mathrm{T}}$ 的股票的均衡价格。
\end{enumerate}
\end{problem}
\begin{solution}
当风险证券支付改为 $(C_\alpha,C_b)$ 时，只需将 11.11 的支付矩阵第二列替换，完全性条件变为 $C_\alpha\ne C_b$，状态价格和二叉树风险中性概率由新的两支付重新解出。若效用改为 CARA，个人一阶条件为 $e^{-ac_{i0}}$ 与状态边际效用的比例关系；Pareto 风险分担下绝对风险容忍为常数，因而每个人承担总风险的固定金额比例。代表性经济人的效用仍为 CARA，SDF 为 $m\propto e^{-aC_1}$，两证券和期权均满足 $P=E[mX]$，期权价格由其两状态支付代入即可。
\end{solution}
\begin{solution}
由 Pareto 一阶条件，$\lambda_i u_i'(f_i(C))$ 在经济人之间相等。对总资源 $C$ 求导并使用隐函数定理，得
\[
f_i'(C)=\frac{T_i(f_i(C))}{\sum_kT_k(f_k(C))}。
\]
因此 $f_i'$ 为常数当且仅当各人的风险容忍在最优分配上相等；这正是风险分担规则线性的充要条件。
\end{solution}
\begin{solution}
\begin{enumerate}[label=(\alph*)]
\item 设状态价格为 $q_s$，个人 $i$ 的 CRRA 参数为 $\gamma_i$。一阶条件为 $p_sc_{is}^{-\gamma_i}=\mu_iq_s$，故 $c_{is}=(p_s/\mu_iq_s)^{1/\gamma_i}$，再用个人财富预算确定 $\mu_i$。
\item 市场出清要求 $c_{1s}+c_{2s}=e_{1s}+e_{2s}$，总资源为 $(100,100,50)$。将两人的需求代入三项资源约束并结合预算约束即可解出正状态价格；相对价格满足 $q_s\propto p_s/[\lambda_1c_{1s}^{-1}+\lambda_2c_{2s}^{-2}]$。
\item 支付向量 $(200,100,50)^{\mathsf T}$ 的均衡价格为 $P=200q_a+100q_b+50q_c$。
\end{enumerate}
\end{solution}
\begin{problem}
\problemno{11.8} 给定 Pareto 最优配置及其最优风险分担规则 $f_i$，令 $T_i(z)=-u_i'(z)/u_i''(z)$ 为经济人 $i$ 的风险容忍函数。
\begin{enumerate}[label=(\alph*)]
    \item 证明
    \[
        f_i'(z)=\frac{T_i(f_i(z))}{\sum_k T_k(f_k(z))}.
    \]
    \item $f_i$ 为线性的充要条件是：$T_i(f_i(z))=T_k(f_k(z))$，$\forall i,k,\forall z$。
\end{enumerate}
\end{problem}
\begin{solution}
由定义 $u_0(C)=\sum_i\lambda_i u_i(f_i(C))$。每个 $u_i$ 单调递增且凹，且 $f_i$ 是最优风险分担规则，故 $u_0'(C)=\lambda_i u_i'(f_i(C))>0$；二阶导数由风险容忍公式为 $u_0''(C)=-u_0'(C)/T_0(C)<0$，所以 $u_0$ 严格递增且凹。对任一总消费过程，先按 $f_i$ 分配可使加权个人效用之和达到最大，因此 $u_0(C_0)+E[u_0(C_1)]$ 正好表示代表性经济人的跨期效用。
\end{solution}
\begin{problem}
\problemno{11.9} 对给定权重 $\lambda_i$，设 $f_i$ 为经济人 $i$ 的 Pareto 最优风险分担规则。定义
\[
    u_0(z)=\sum_i \lambda_i u_i(f_i(z)).
\]
证明：
\begin{enumerate}[label=(\alph*)]
    \item $u_0$ 是严格单调递增且凹的；
    \item $u_0(C_0)+E[u_0(C_1)]$ 可表示代表性经济人的效用(在适当权重下)。
\end{enumerate}
\end{problem}
\begin{problem}
\problemno{11.10} 证明：如果每个经济人是绝对风险厌恶递减的(DARA)，那么代表性经济人也是绝对风险厌恶递减的。
\end{problem}
\begin{solution}
由 $T_0(C)=\sum_iT_i(f_i(C))$ 及 $f_i'(C)\ge0$，每个 DARA 经济人的风险容忍 $T_i$ 随财富增加而增加。因此
\[
T_0'(C)=\sum_iT_i'(f_i(C))f_i'(C)\ge0。
\]
代表性经济人的风险容忍递增，等价于 $A_0(C)=1/T_0(C)$ 递减，所以代表性经济人也是 DARA。
\end{solution}
\begin{problem}
\problemno{11.11} 两时点单期情形下，假设期末有两个状态 $\alpha$ 和 $b$，它们发生的概率分别为 $\pi$ 和 $1-\pi$。市场上有两只证券：证券 1 为无风险证券，证券 2 为风险证券，支付向量为 $(uC,dC)$，其中 $C>0$，$0<d<u$。经济中有 $K$ 个经济人，经济人 $k$ 的禀赋是 0 期 $s_k$ 个单位的消费品和 $s_k$ 单位的证券 2，其中 $s_k>0$ 且 $\sum_k s_k=1$。所有经济人的效用函数为
\[
    U(c)=\ln c_{k,0}+\rho\bigl[\pi\ln c_{k,\alpha}+(1-\pi)\ln c_{k,b}\bigr],\qquad \rho>0.
\]
\begin{enumerate}[label=(\alph*)]
    \item 证明证券市场是完全的；
    \item 求解每个经济人的最优消费-投资选择；
    \item 求两个证券的均衡解。
\end{enumerate}
\end{problem}
\begin{solution}
\begin{enumerate}[label=(\alph*)]
\item 支付矩阵为 $\begin{bmatrix}1&uC\\1&dC\end{bmatrix}$，行列式 $C(d-u)\ne0$，故两证券张成两个状态，市场完全。
\item 令证券 1 持有量为 $b_k$、证券 2 持有量为 $\theta_k$。预算约束为 $c_{k0}+q_1b_k+q_2\theta_k=s_kq_2$，状态消费为 $c_{k\alpha}=b_k+uC\theta_k$、$c_{kb}=b_k+dC\theta_k$。对数效用的一阶条件给出 $c_{k\alpha}/c_{kb}$ 等于状态价格调整后的概率比，代入预算约束即可得到最优 $b_k,\theta_k$。
\item 市场出清为 $\sum_kb_k=0$（无风险证券净供给为零）、$\sum_k\theta_k=1$（风险证券总供给为一），以及两状态消费分别等于总支付 $uC,dC$。联立出清条件和一阶条件得到两证券均衡价格，且状态价格严格为正。
\end{enumerate}
\end{solution}
\begin{problem}
\problemno{11.12} 在第 11.11 题的框架下，将价格过程与期权定价中二叉树模型进行比较，可以得出什么结论？
\end{problem}
\begin{solution}
两种模型都由两个状态支付和无套利条件决定状态价格。均衡中的证券价格等于状态价格加权支付，二叉树中的风险中性概率只是将状态价格除以无风险债券价格后的归一化形式。因此二叉树期权定价可以看作完全市场均衡定价在两状态情形下的递推表达。
\end{solution}
\begin{problem}
\problemno{11.13} 在第 11.11 题的框架下：
\begin{enumerate}[label=(\alph*)]
    \item 构建代表性经济人；
    \item 计算代表性经济人的效用函数中每个经济人的权重。它依赖什么？为什么？
    \item 证明代表性经济人的偏好与禀赋在各经济人中的分布无关；
    \item 用基于消费的 CAPM 给两证券定价；
    \item 计算风险资产的股权溢价；
    \item 利率和股票形成的股权风险溢价依赖于什么？解释它们对当前消费水平、消费增长率的期望值 $E[C_1/C_0]$ 和波动率 $\operatorname{var}(C_1/C_0)$ 以及时间偏好率 $\rho$ 的依赖性；
    \item 求以证券 2 为标的资产、执行价格为 $K$ 的看涨期权的价格，并证明基于消费的 CAPM 也适用于期权。
\end{enumerate}
\end{problem}
\begin{solution}
\begin{enumerate}[label=(\alph*)]
\item 由于所有个人具有相同的对数偏好，Pareto 加总后的代表性经济人仍为对数效用：$U_0=\ln C_0+\rho E[\ln C_1]$。
\item 代表性权重等于个人预算约束中的财富权重，取决于初始禀赋和 Pareto 权重；禀赋改变财富分配但不改变共同的边际替代率形式。
\item 风险分担使每个人状态消费是总消费的固定比例，因此代表性偏好只依赖总消费过程，不依赖总禀赋在个人之间如何分配。
\item SDF 为 $m=\rho C_0/C_1$（状态逐一理解），两证券价格满足 $P_j=E[mX_j]$。
\item 风险资产股权溢价为 $E[R_2]-R_f=-R_f\,\operatorname{Cov}(m,R_2)/E[m]$，由消费增长与资产支付的协方差决定。
\item 无风险利率取决于 $E[C_1/C_0]$、其波动率和时间偏好 $\rho$；消费增长越高通常降低贴现、提高无风险收益，波动越大且经济人越风险厌恶则风险溢价越高。
\item 看涨期权支付为 $(X_2-K)^+$，其价格为 $E[m(X_2-K)^+]$。因为该支付也是资产支付，故同一基于消费的 CAPM 定价关系自动适用于期权。
\end{enumerate}
\end{solution}
\begin{problem}
\problemno{11.14} 修改第 11.11 题和第 11.13 题中的两个条件并再次求解第 11.11--11.13 题中各小问：
\begin{enumerate}[label=(\alph*)]
    \item 风险资产的支付向量修改为 $(C_\alpha,C_b)$；
    \item 效用函数修改为 CARA 偏好，即 $u(z)=-e^{-az}$。
\end{enumerate}
\end{problem}


% !TeX root = ../../金融经济学_熊和平_习题.tex
% 源：113071.pdf p.10--11（印张约 273；本章小结在扫描件前页有残缺）

\zhchapter{基于消费的资产定价模型}

\begin{tcolorbox}[summary,title=本章小结]
基于消费的资产定价模型在完全市场框架下给出了资产价格与消费之间的简洁关系，并可在一定条件下得到随机贴现因子、无风险利率等的显示解。然而，这些简洁的定价方程在实证中遇到挑战，形成了股权溢价之谜、无风险利率之谜和消费平滑之谜等。对这些谜题的探索推动了金融经济学的发展。
\end{tcolorbox}

\xhpconceptitems{
    \item 随机贴现因子（SDF），股权溢价
    \item 股权溢价之谜，无风险利率之谜，股价波动之谜
    \item 递归偏好，跨时替代弹性
}

% !TeX root = ../../金融经济学_熊和平_习题.tex

\begin{problem}
\problemno{12.1} 随机贴现因子定价模型(SDF)为什么也叫基于消费的资产定价模型？
\end{problem}
\begin{problem}
\problemno{12.2} 本章的随机贴现因子与上一章的随机贴现因子是否相同？
\end{problem}
\begin{problem}
\problemno{12.3} 为什么说 SDF 模型具有普适性？
\end{problem}
\begin{problem}
\problemno{12.4} 基于 SDF 的模型与 CAPM 的联系与区别是什么？
\end{problem}
\begin{problem}
\problemno{12.5} 什么是股权溢价之谜？
\end{problem}
\begin{problem}
\problemno{12.6} 什么是无风险利率之谜？
\end{problem}

\begin{problem}
\problemno{12.7} 两时点单期情形下假设期末有 $S$ 种状态，则基于消费的定价模型表明任何资产的价格 $P_i$ 和期末的支付 $X_i$ 满足
\[
    P_i=\sum_{s=1}^{S}\pi_s m_s X_{is}=E[mX_i].
\]
如果经济中一个投资者的效用函数为 $u(c_0,c_1)$，$\omega\in\Omega$ 为均衡状态下的消费水平，试用投资者的均衡消费水平表示状态价格。
\end{problem}
\begin{problem}
\problemno{12.8} 假设期末有两种状态，市场上有两种证券：一个是无风险证券，其相对收益为 $1.05$；另一个是风险证券，其期初价格为 $6$，期末支付向量为 $(10,5)$。
\begin{enumerate}[label=(\alph*)]
    \item 计算两个状态价格和各状态的风险中性概率；
    \item 若期末两状态等概率发生，计算两个状态所对应的随机贴现因子。
\end{enumerate}
\end{problem}
\begin{problem}
\problemno{12.9} 两时点、两状态单期情形下，一份期权合约在期末状态 1 下的支付为 $3$，状态 2 下的支付为 $0$，0 期市场价格为 $1$。一份相同标的资产的远期合约在状态 1 下的支付为 $3$，状态 2 下的支付为 $-2$。试求无风险利率及两状态的风险中性概率。
\end{problem}
\begin{problem}
\problemno{12.10} 假如随机贴现因子定价关系与 CAPM 都成立，其中 CAPM 模型表述为 $E[R_i]=R_f+\beta_i\gamma$，随机贴现因子表述为 $m=a+bR_m$，$R_m$ 是市场投资组合的收益，其方差记为 $\sigma_m^2$。
\begin{enumerate}[label=(\alph*)]
    \item 将 $\gamma$ 表示为 $a$、$b$、$E[R_m]$ 和 $\sigma_m^2$ 的函数(可从 $E[m(R_i-R_f)]=0$ 入手)；
    \item 将 $a$ 和 $b$ 表示为 $R_f$、$E[R_m]$ 和 $\sigma_m^2$ 的函数形式(可从 $E[mR_f]=1$ 入手)；
    \item 利用 (2) 中的结果证明 $\gamma=E[R_m]-R_f$。
\end{enumerate}
\end{problem}
\begin{problem}
\problemno{12.11} 假设存在 $N$ 个风险资产，其期望收益向量为 $\mu$，协方差矩阵为 $\Sigma$。记 $\mu_m=E[m]$ 和 $\sigma_m^2$ 分别为贴现因子的均值和方差。
\begin{enumerate}[label=(\alph*)]
    \item 定义 $m_0=\mu_m+(1-\mu_m\mu)^{\mathrm{T}}\Sigma^{-1}(R-\mu)$，其中 $1$ 为 $N$ 维单位向量，$R$ 为 $N$ 个资产的收益向量。请证明：$1=E[m_0 R]$，即 $m_0$ 也是一个贴现因子。
    \item 利用柯西-施瓦兹不等式证明
    \[
        \sigma_m^2\ge \operatorname{var}[m]=(1-\mu_m\mu)^{\mathrm{T}}\Sigma^{-1}(1-\mu_m\mu).
    \]
\end{enumerate}
\end{problem}


% !TeX root = ../../金融经济学_熊和平_习题.tex
% 源：微信图片_20260928151417_3591_382.jpg（教材印张 293）

\zhchapter{生产性经济与资源配置}

% !TeX root = ../../金融经济学_熊和平_习题.tex
% 源：微信图片_20260928151417_3591_382.jpg（教材印张 293）

\begin{problem}
\problemno{13.1} 在纯交换经济基础上，有哪些资源配置的一般拓展路径？
\end{problem}
\begin{problem}
\problemno{13.2} 储蓄或储存一般可以给消费可行集带来怎样的影响？
\end{problem}
\begin{problem}
\problemno{13.3} 生产技术通常如何影响资源配置？
\end{problem}
\begin{problem}
\problemno{13.4} 简单说明纯交换经济与生产性经济的区别与联系。
\end{problem}

\begin{problem}
\problemno{13.5} 经济中有 1 个经济人，他具有对数效用函数。期末有等概率的两个状态：状态 $a$ 和状态 $b$。经济人的初始禀赋为期初 \num{200} 个消费品，但他可以通过储存将消费品跨时转移。
\begin{enumerate}[label=(\alph*)]
    \item 当利率为 $1.05\%$ 时，分别确定时间偏好率为 $1$ 和 $0.5$ 时的最优资源配置；
    \item 当利率为 $-5\%$ 时，分别确定时间偏好率为 $1$ 和 $0.5$ 时的最优资源配置；
    \item 进一步如果该经济人期末也有禀赋，分别为 $a$ 状态下 \num{100} 个消费品、$b$ 状态下 \num{50} 个消费品，分别求解 (1)(2) 中的问题。
\end{enumerate}
\end{problem}
\begin{problem}
\problemno{13.6} 经济中有 1 个经济人，他具有对数效用函数。期末有等概率的两个状态：状态 $a$ 和状态 $b$。经济人的初始禀赋为期初 \num{200} 个消费品，他可以如下生产技术进行资源跨时转移：
\[
    y_1=f_{1,\omega}(I)=
    \begin{cases}
        \sqrt{I}, & \omega=a,\\
        2\sqrt{I}, & \omega=b.
    \end{cases}
\]
\begin{enumerate}[label=(\alph*)]
    \item 求最优投资数量；
    \item 求状态价格和最优消费计划。
\end{enumerate}
\end{problem}
\begin{problem}
\problemno{13.7} 经济中有 1 个经济人，他具有对数效用函数。期末两个状态 $a$ 和 $b$ 的概率分别为 $\pi_a$ 和 $\pi_b$。经济人的初始禀赋为期初 \num{200} 个消费品，他可以如下生产技术进行资源跨时转移：
\[
    y_1=f_{1,\omega}(I)=
    \begin{cases}
        \lambda_a I, & \omega=a,\\
        \lambda_b I, & \omega=b.
    \end{cases}
\]
证券市场上存在无风险债券，其利率为 $r_f$。
\begin{enumerate}[label=(\alph*)]
    \item 给出经济人的最优决策，包括投资决策和最优消费；
    \item 求解他的最优投资数量和消费计划；
    \item 分析最优投资与利率的关系；
    \item 假设该经济人没有任何商品禀赋，回答上述 (1) 至 (3) 问。
\end{enumerate}
\end{problem}
\begin{problem}
\problemno{13.8} 在第 3 题的背景下引进经济人 2，经济人 1 只有生产技术禀赋，经济人 2 的禀赋是期初 \num{200} 个消费品：
\begin{enumerate}[label=(\alph*)]
    \item 求两个人的最优消费计划和状态价格；
    \item 求经济人 1 的最优投资。
\end{enumerate}
\end{problem}


% !TeX root = ../../金融经济学_熊和平_习题.tex
% 源：微信图片_20260928151500_3593_382.jpg（教材印张 311）

\zhchapter{异质性投资者}

\begin{tcolorbox}[summary,title=本章小结]
在行为金融中所重现的一些心理特征都可能形成特定的异质性；本章的异质性属于一类最基本的异质。即使是上述基本的异质性也会对资产定价和最终的资源配置产生很大的影响。

学术界针对经济人的异质性通常有三大假设：一是市场平衡假设，认为众多异质性经济人相互作用最终可以将诸多个性“平衡”掉，从而形成貌似只有一个经济人的情形；第二是市场选择假设，借助生物中的“物竞天择、自然选择”思想，认为市场将会淘汰“错误的”经济人，最终使聪明的经济人留下，从而使得资本市场被聪明的经济人所主导；第三是市场学习假设，认为经济人利用可获得的信息更新自己对市场的认知，从而形成正确的预期并最终使市场趋于理想化的市场。三种假设的模型化和实证分析都推动金融经济学理论的发展。

市场平衡假设强调所有经济人的“集体”作用，不同经济人共同作用形成了资本市场的加总理论，异质禀赋下的加总形成加总性质，CRRA 偏好的加总则形成一类简单的、理性化的情形。
\end{tcolorbox}

\begin{tcolorbox}[knowledge,title=重要概念]
异质性投资者，异质禀赋，异质偏好，异质信念，时间偏好，资本市场加总，市场选择，市场平衡，投资者学习，贝叶斯更新，加总性质。
\end{tcolorbox}

% !TeX root = ../../金融经济学_熊和平_习题.tex

\xhpthoughts
\begin{problem}
\problemno{1} 投资者的异质性一般包括哪几个方面？
\end{problem}
\begin{problem}
\problemno{2} 什么是投资者的加总性质？
\end{problem}
\begin{problem}
\problemno{3} 关于异质性经济人有哪三大假设？
\end{problem}
\begin{problem}
\problemno{4} 加总性质与特定效用函数下的加总有哪些区别与联系？
\end{problem}
\begin{problem}
\problemno{5} 如何理解理性学习与贝叶斯更新？
\end{problem}

\xhpexercises
\begin{problem}
\problemno{1} 在例 14.14 中取 $B=1$，$A=0$ 时，效用函数为对数偏好。请在对数偏好下对经济人进行加总。
\end{problem}
\begin{problem}
\problemno{2} 假设经济中两个经济人的时间偏好分别为 $0.5$ 和 $1$，他们都具有 CRRA 效用函数，他们的相对风险厌恶系数为 $\gamma=1$，分别求例 14.1 中三种异质禀赋下的状态价格和资源配置。
\end{problem}
\begin{problem}
\problemno{3} 在上题的条件下进一步假设两个经济人具有异质性偏好，相对风险厌恶系数分别为 $1$ 和 $0.1$。在此条件下，分别求例 14.1 中三种异质禀赋下的状态价格和资源配置。
\end{problem}
\begin{problem}
\problemno{4} 假设经济中两个经济人的时间偏好均为 $1$，他们都具有 CRRA 效用函数，他们的相对风险厌恶系数为 $1$ 和 $0.5$，分别求例 14.1 中三种异质禀赋下的状态价格和资源配置。
\end{problem}


% !TeX root = ../../金融经济学_熊和平_习题.tex
% 源：113071.pdf p.14（印张 330）；微信 51600--51601（印张 330--331）

\zhchapter{家庭金融理论}

\begin{tcolorbox}[summary,title=本章小结]
随着经济的发展，家庭金融的规模越来越大。在所有发达国家，家庭使用的金融服务和产品在金融业中占据了相当大的比例。在某种程度上，市场规模是衡量重要性的一个指标，家庭的财务状况至少应该得到与企业财务状况同样多的关注。另一方面，家庭决策及其结果往往取决于作出决策的体制环境，同时还要面临复杂的金融市场和一些具体的监管干预措施。因此，家庭金融有着与传统金融经济学不同的研究目标，家庭金融面临的独特挑战之一是研究金融复杂性及其对家庭决策和福利的影响。

研究家庭金融首先要从家庭的禀赋和约束出发。家庭的资本禀赋包括人力资本和有形财富。人力资本是很多家庭一生中的主要收入来源，它是家庭在剩余生命周期中各年预期的来自劳动力收益的折现值。人力资本是通过正规教育或工作经验慢慢积累起来的，其价值很难评估，而且不可交易，属于家庭的背景风险。有形财富则呈现多样性，不同家庭的有形财富差异巨大。因此，可以说家庭的异质性将使金融经济学从纯理论走到脚踏实地服务于一个个具体的家庭。

家庭的投资决策从理论上讨论家庭应该如何投资。家庭的生命期长久，而且家庭在不同时期投资的目标有所不同，因此一般将家庭划分为不同的期限，从而形成基于生命周期的模型。家庭处在不同生命周期的储蓄行为形成基于生命周期的储蓄理论；不同时期的投资组合决策则形成基于生命周期的组合理论。莫顿—莫森—萨缪尔森模型是一类关于家庭一生投资组合的基准模型。由于家庭异质性显著，家庭的基于生命周期的组合理论不断完善和发展，相应的模型也更加复杂。

比照理论模型的结论，实证家庭金融表明家庭在金融决策中存在一些错误。比如，金融学家们列举了四个投资错误：很多家庭不参与资本市场，从而失去了利用股票市场提高家庭福利的机会；参与股票市场的家庭往往没有对投资资金进行合理分散化投资，在承担相同风险的条件下没有获得最大的收益；风险分担过于简单，使家庭暴露于过多不必要的风险；在购买住房做按揭贷款时，抵押再融资模式不科学。通过家庭金融教育和家庭金融工程可以对这些投资错误在一定程度上进行纠正。
\end{tcolorbox}

\xhpconceptitems{
    \item 家庭金融，消费金融，生命周期的储蓄模型，家庭生命周期
    \item 不确定性，模糊性，MMS 模型
    \item 人力资产，有形财富，投资错误，家庭金融工程
}

% !TeX root = ../../金融经济学_熊和平_习题.tex

\begin{problem}
\problemno{15.1} 什么是家庭金融？如何理解家庭金融从金融学中分离出来是必要的？
\end{problem}
\begin{solution}
家庭金融研究家庭在储蓄、借贷、住房、保险、退休和投资中的决策及其与金融市场的互动。家庭同时拥有劳动收入、人力资本、住房和非交易资产，并面临生命周期、家庭成员和制度约束，因此不能直接套用只有可交易金融资产的代表性投资者模型，独立研究有助于解释真实资产配置和福利差异。
\end{solution}
\begin{problem}
\problemno{15.2} 简述家庭金融理论与异质性投资者理论的关系，家庭作为经济人有哪些异质性？
\end{problem}
\begin{solution}
家庭金融把异质性投资者框架应用于家庭决策。家庭在财富、年龄、劳动收入风险、风险偏好、金融知识、健康状况、婚姻结构、税收和借贷约束上都可能不同，这些异质性造成储蓄率、住房和股票参与率差异。
\end{solution}
\begin{problem}
\problemno{15.3} 为什么说家庭资源禀赋和面临的经济环境更加复杂？
\end{problem}
\begin{solution}
家庭资源不仅有金融财富，还包括劳动收入、人力资本、住房和耐用品；支出还受子女、教育、医疗和养老需求影响。收入和支出风险具有长期性、不可交易性和相关性，且家庭受到税制、社会保险、信贷限额和信息成本约束，所以其预算集远比单一证券投资者复杂。
\end{solution}
\begin{problem}
\problemno{15.4} 什么是模糊性？简述风险、不确定性和模糊性的区别与联系。
\end{problem}
\begin{solution}
风险是结果集合和概率都已知；不确定性是结果可能已知但客观概率无法确定；模糊性强调对概率分布本身存在多种合理判断。三者都涉及未来状态，风险是最可量化的情形，模糊性可看作对概率的二阶不确定，通常会产生额外的模糊厌恶。
\end{solution}
\begin{problem}
\problemno{15.5} 家庭投资决策有哪些错误？如何理解这些错误？
\end{problem}
\begin{solution}
教材归纳的四类典型错误是：不参与资本市场从而放弃股票市场的福利改进；参与后分散不足、在相同风险下收益偏低；风险分担安排过于简单、承担不必要风险；住房按揭与再融资模式不科学（如高息负债与低息存款并存）。理解这些错误需结合有限理性、金融知识、社会规范与产品复杂性，而非仅归因于“不理性”。
\end{solution}
\begin{problem}
\problemno{15.6} 什么是基于生命周期的投资模型？
\end{problem}
\begin{solution}
生命周期模型把家庭的劳动收入、消费、储蓄和投资放在多个年龄阶段联合决定。年轻时人力资本高、金融财富少，通常积累资产；中年收入和财富较高，进行风险分散；退休后依靠金融资产和养老金消费。最优股票比例取决于剩余寿命、劳动收入风险、风险厌恶和退休制度，而非固定不变。
\end{solution}
\begin{problem}
\problemno{15.7} 某家庭为了购买住房在银行办理了 \num{50} 万元按揭贷款，贷款利率为 $6\%$，同时他们在其他银行拥有 \num{50} 万元定期存款，存款利率为 $3\%$。当问及为何购买住房使用贷款时，答曰：他熟悉的朋友在买房时都使用了贷款，何况银行给他们的贷款利率比一般的贷款利率优惠。请分析该家庭该项金融决策的错误。
\end{problem}
\begin{solution}
该家庭同时借款和持有同额存款，相当于以 $6\%$ 的利率借入资金再以 $3\%$ 的利率投资，忽略税收和流动性差异时每年确定损失为 $500000(6\%-3\%)=15000$ 元。决策依据“朋友都这样做”是从众和熟悉性偏差，优惠贷款利率低于一般贷款并不能说明低于其存款机会成本。正确比较应考虑提前还贷收益、存款流动性、税收、违约风险和住房的非金融收益；若无特殊约束，先用存款偿还高利率贷款更有效率。
\end{solution}



\appendix
% !TeX root = ../../金融经济学_熊和平_习题.tex
% 源：微信图片_20260928151601_3596_382.jpg（教材印张 331 起；性质 2 在扫描件末处截断）

\chapter{布朗运动及其简单描述}

\begin{tcolorbox}[summary,title=附录提要]
通常假设股票收益率由两部分确定：确定性的部分和随机性的部分，随机性由维纳过程或布朗运动刻画。参照 John Hull 的描述，称随机过程 $z$ 服从维纳过程，若 $z$ 具有如下两个性质。

\textbf{性质 1}：在很短的时间间隔 $\Delta t$ 内的增量 $\Delta z$ 可以表示为
\[
    \Delta z=\epsilon\sqrt{\Delta t}
    \tag{A.1}
\]
其中 $\epsilon$ 服从标准正态分布，即 $\epsilon\sim N(0,1)$；取 $\Delta t\to 0$ 时，(A.1) 可记为
\[
    \mathrm{d}z=\epsilon\sqrt{\mathrm{d}t}.
    \tag{A.2}
\]

\textbf{性质 2}：$z$ 在任何两个不相互重叠的时间区间上的增量相互独立。% 原书后续文字在扫描件末处截断，待补扫。
\end{tcolorbox}


\end{document}
